%%%%%%%%%%%%%%%%%%%%%%%%%%%%%%%%%%%%%%%%%%%%%%%%%%%%%%%%%%%%%%%%%%%%%%
% Lớp 11 | Các số đặc trưng đo xu thế trung tâm cho mẫu số liệu | Số trung bình và mốt của mẫu số liệu ghép nhóm
% Dạng 1D5?1-4: Mốt
% Số câu: 117  (sinh tự động bởi tach_cau.py; từ nay thầy sửa trực tiếp file này)
%%%%%%%%%%%%%%%%%%%%%%%%%%%%%%%%%%%%%%%%%%%%%%%%%%%%%%%%%%%%%%%%%%%%%%

%%%%%%%%%%%%%%%% TRẮC NGHIỆM 4 PHƯƠNG ÁN %%%%%%%%%%%%%%%%

% nguồn: Ok_Toan11_TracNghiem_Lan2
\begin{ex}%[1D5N1-4]%[ma:fb95bfca]
	Mức thưởng tết cho các nhân viên của một công ty được thống kê  trong bảng sau:
	\begin{center}
		\begin{tabular}{|c|c|c|c|c|c|}
			\hline 
			Mức thưởng tết &  $[5;10)$ & $[10;15)$ & $[15;20)$ & $[20;25)$ & $[25;30)$\\
			 \hline
			Số nhân viên & $13$ & $35$ & $47$ & $25$ & $10$\\ 
			 \hline
		\end{tabular}
	\end{center}
	Số mốt của mẫu số liệu ghép nhóm trên  là
	 \choice
		{$0$} 
		{\True  $1$}
		{$2$} 
		{$3$} 
	\loigiai{ 
		Tần số lớn nhất của mẫu  số liệu ghép nhóm trên là $47$.\\
		Nên nhóm chứa mốt của mẫu số liệu trên là  nhóm $[15;20)$.\\
		Do đó mẫu số liệu này chỉ có $1$ mốt. 
	 }
\end{ex}

% nguồn: Ok_Toan11_TracNghiem_Lan2
\begin{ex}%[1D5N1-4]%[ma:6b5bff6b]
	Gọi $i$  là nhóm có tần số lớn nhất. Gọi $u$, $g$, ${n_i}$ lần lượt là đầu mút trái, độ dài và tần số của nhóm $i$; $n_{i-1}$, $n_{i+1}$ lần lượt là tần số của nhóm $i-1$, nhóm $i+1$. Gọi $M_o$ là Mốt của mẫu số liệu ghép nhóm, mệnh đề nào sau đây đúng?
	 \choice
		{\True $M_o=u+(\dfrac{n_i-n_{i-1}}{2n_i-n_{i-1}-n_{i+1}})\cdot  g$}
		{$M_o=u+(\dfrac{n_i\cdot  n_{i-1}}{2n_i-n_{i+1}-n_{i-1}})\cdot g$}
		 {$M_o=u\cdot (\dfrac{n_i-n_{i-1}}{2n_i-n_{i-1}-n_{i+1}})-g$}
		{$M_o=u+(\dfrac{n_{i+1}-n_{i-1}}{2n_i-n_{i-1}-n_{i+1}})\cdot  g$}
	 \loigiai{
		 Theo lí thuyết ta có $M_o=u+(\dfrac{n_i-n_{i-1}}{2n_i-n_{i-1}-n_{i+1}})\cdot g$.
	 }
\end{ex}

% nguồn: Ok_Toan11_TracNghiem_Lan2
\begin{ex}%[1D5N1-4]%[ma:20276699]
	Thời gian truy cập internet mỗi buổi tối của một số học  sinh được cho trong bảng sau
	\begin{center}
		\begin{tabular}{|c|c|c|c|c|c|}
			\hline 
			Thời gian & $[9{,}5;12{,}5)$  & $[12{,}5;15{,}5)$ & $[15{,}5;18{,}5)$ & $[18{,}5;21{,}5)$ & $[21{,}5;24{,}5)$\\
			 \hline
			Số học sinh & $3$ & $12$ & $15$ & $24$ & $2$\\ 
			\hline 
		\end{tabular}
	\end{center}
	Nhóm chứa  mốt là nhóm nào?
	\choice 
		 {$[9{,}5;12{,}5)$}
		 {$[12{,}5;15{,}5)$}
		{$[15{,}5;18{,}5)$} 
		 {\True $[18{,}5;21{,}5)$}
	 \loigiai{
		Vì tần số lớn nhất là $24$ nên  nhóm chứa mốt là $[18{,}5;21{,}5)$.
	} 
\end{ex}

% nguồn: Ok_Toan11_TracNghiem_Lan2
\begin{ex}%[1D5N1-4]%[ma:08c44944]
	Thời gian (phút) truy cập internet  mỗi buổi tối của một số học sinh được cho trong bảng sau
	\begin{center}
		\begin{tabular}{|c|c|c|c|c|c|}
			\hline 
			Thời gian (phút) & $[9{,}5;12{,}5)$ &  $[12{,}5;15{,}5)$ & $[15{,}5;18{,}5)$ & $[18{,}5;21{,}5)$ & $[21{,}5;24{,}5)$ \\
			 \hline
			Số học sinh & $3$ & $12$ &  $15$ & $24$ & $2$ \\
			 \hline
		\end{tabular}
	\end{center}
	 Nhóm chứa mốt là nhóm nào?
	 \choice
		 {$\left[9{,}5;12{,}5 \right)$}
		{$\left[12{,}5;15{,}5  \right)$}
		{$\left[15{,}5;18{,}5  \right)$}
		{\True $\left[18{,}5;21{,}5  \right)$}
	 \loigiai{
		Vì tần số lớn nhất là $24$ nên nhóm chứa mốt là  $\left[18{,}5;21{,}5 \right)$.
	 }
\end{ex}

% nguồn: Ok_Toan11_TracNghiem_Lan2
\begin{ex}%[1D5N1-4]%[ma:189bf738]
	Khảo sát thời gian tự học bài ở nhà của một số em học sinh lớp $11$ thu được mẫu ghép nhóm số lượng  như sau:
	\begin{center}
		\begin{tabular}{|c|c|c|c|c|c|}
			\hline 
			Thời  gian(phút) & $[0;30)$ & $[30;60)$ & $[60;90)$ & $[90;120)$ & $[120;150)$\\
			\hline 
			Số  học sinh & $9$ & $10$ & $9$ & $15$ & $7$\\
			\hline 
		\end{tabular}
	\end{center}
	 Nhóm chứa mốt là
	 \choice
		{$[30;60)$} 
		{$[60;90)$} 
		{\True  $[90;120)$}
		 {$[120;150)$}
	\loigiai{ 
		Nhóm $[90;120)$ có tần số cao nhất  là $15$.
	} 
\end{ex}

% nguồn: Ok_Toan11_TracNghiem_Lan2
\begin{ex}%[1D5N1-4]%[ma:cc6619a9]
	 Nhiệt độ trung bình của $12$ tháng tại thành phố Hà Nội từ năm $1961$ đến hết năm $1990$ ($30$ năm) được cho trong bảng sau
	\begin{center}
		\begin{tabular}{|l|c|c|c|c|c|}
			\hline Nhiệt độ ($^\circ$C) & $[12;14)$ & $ [14;16)$ & $[16;18) $ & {$[18;20)$} &  {$[20;22)$} \\
			\hline Tần số &  $1$ & $3$ & $12$ & $9$& $5$ \\
			\hline 
		\end{tabular}
	\end{center}
	mốt của  mẫu số liệu trên là
	\choice 
		 { $15{,}7$}
		{  $16{,}0$}
		{ \True  $17{,}5$}
		{ $17{,}0$} 
	 \loigiai{
		Nhóm chứa mốt của mẫu số liệu  trên là nhóm $[16;18)$.\\
		 Nên ta có $u_m=16,n_{m-1}=3,n_m=12,n_{m+1}=9,u_{m+1}-u_m=18-16=2.$\\
	mốt của mẫu số liệu trên là $M_o=16+\dfrac{12-3}{(12-3)+(12-9t)}\cdot 2=17{,}5.$} 
\end{ex}

% nguồn: Ok_Toan11_TracNghiem_Lan2
\begin{ex}%[1D5N1-4]%[ma:1c71ad06]
	Khảo sát thời gian tập thể dục trong ngày của một số học sinh khối $11$ thu được  mẫu số liệu ghép nhóm sau:
	\begin{center}
		\begin{tabular}{|c|c|c|c|c|c|}
			 \hline
			Thời gian (phút) & $[0;20)$  & $[20;40)$ & $[40;60)$ & $[60;80)$ & $[80;100)$\\
			\hline 
			Số học sinh & $5$ &  $9$ & $12$ & $10$ & $6$\\
			\hline 
		\end{tabular}
	\end{center}
	 Mẫu số liệu ghép nhóm này có số mốt là
	 \choice
		 {$0$}
		{\True  $1$}
		 {$2$}
		 {$3$}
	 \loigiai{
		Vì đây là mẫu số liệu ghép nhóm  và tần số các nhóm khác nhau nên có $1$ mốt.
	} 
\end{ex}

% nguồn: Ok_Toan11_TracNghiem_Lan2
\begin{ex}%[1D5N1-4]%[ma:9cdbafc8]
	Khi thống kê chiều cao của học sinh khối $12$ trong một trường trung học, ta thu được mẫu số liệu ghép nhóm sau 
	\begin{center}
		\begin{tabular}{|c|c|c|c|c|c|c|}
			\hline 
			Chiều cao &  $[150;156)$ & $[156;162)$ & $[162;168)$ & $[168;174)$ & $[174;180)$ & $[180;186)$\\
			 \hline
			Số học sinh & $5$ & $18$ & $40$ & $26$  & $8$ & $3$\\
			 \hline
		\end{tabular}
	\end{center}
	Nhóm chứa mốt của mẫu số liệu ghép nhóm  trên là
	\choice 
		{$[150;156)$} 
		{\True  $[162;168)$}
		 {$[168;174)$}
		 {$[180;186)$}
	\loigiai{ 
		Nhóm chứa mốt của mẫu số liệu ghép nhóm là nhóm có tần số lớn  nhất.\\
	Do đó  nhóm chứa mốt của mẫu số liệu ghép nhóm trên là nhóm $[162;168)$.}
\end{ex}

% nguồn: Ok_Toan11_TracNghiem_Lan2
\begin{ex}%[1D5N1-4]%[ma:459ccf0e]
	Khảo sát thời gian tập thể dục của một số  học sinh khối $11$ thu được mẫu số liệu ghép nhóm sau:
	\begin{center}
		\begin{tabular}{|c|c|c|c|c|c|}
			\hline 
			Thời gian (phút) & $[0;20)$ & $[20;40)$ & $[40;60)$  & $[60;80)$ & $[80;100)$\\
			 \hline
			Số học  sinh & $5$ & $9$ & $12$ & $10$ & $6$\\
			 \hline
		\end{tabular}
	\end{center}
	 Nhóm chứa mốt của mẫu số liệu trên là
	\choice 
		 {$[20;40)$}
		 {$[60;80)$}
		{\True $[40;60)$} 
		{$[80;100)$} 
	\loigiai{ 
		Mốt $M_o$ chứa trong nhóm $[40;60)$ vì nhóm này  có tần số lớn nhất.
	} 
\end{ex}

% nguồn: Ok_TracNghiemLop11_Lan1
\begin{ex}%[1D5N1-4]%[ma:2fc8b0d5]
	Kết quả khảo sát cân nặng của $25$ quả cam ở lô hàng $A$  được cho ở bảng sau
	\begin{center}
		{\tabcolsep = 2mm 
		\begin{tabular}{|c|c|c|c|c|c|}
			 \hline
			Cân nặng & $[150;  155)$ & $[155; 160)$ & $[160; 165)$ & $[165; 170)$ & $[170; 175)$ \\
			\hline 
			Số quả cam ở lô hàng $A$ & $1$ &  $3$ & $7$ & $10$ & $4$\\
			 \hline
		\end{tabular}
		 }
	\end{center}
	Nhóm chứa mốt  là nhóm nào?
	\choice 
		{$[150;  155)$}
		 {$[155; 160)$}
		 {\True $[165; 170)$}
		{$[170;  175)$}
	\loigiai{ 
		Nhóm chứa mốt của mẫu số liệu ghép nhóm là nhóm có tần số  lớn nhất.\\
		Nhóm chứa  mốt của mẫu số liệu ghép nhóm là nhóm $[165; 170)$.
	 }
\end{ex}

% nguồn: Ok_TracNghiemLop11_Lan1
\begin{ex}%[1D5N1-4]%[ma:2aeea2e4]
	Khảo sát thời gian chạy bộ trong một ngày của một số học sinh khối $11$ thu được mẫu số liệu ghép nhóm sau 
	\begin{center}
		 {\tabcolsep = 2mm
		\begin{tabular}{|c|c|c|c|c|c|}
			 \hline
			Thời gian (phút) & $[0; 20)$  & $[20; 40)$ & $[40; 60)$ & $[60; 80)$ & $[80; 100)$ \\
			 \hline
			Số  học sinh & $5$ & $9$ & $12$ & $10$ & $6$\\
			 \hline
		\end{tabular}
		 }
	\end{center}
	Nhóm chứa mốt của mẫu số liệu ghép nhóm  trên là
	 \choice
		{$[20;  40)$}
		{\True $[40;  60)$}
		 {$[60; 80)$}
		{$[80; 100)$} 
	 \loigiai{
		Mốt  $\mathrm{M}_0$ chứa trong nhóm $[40;60)$.
	} 
\end{ex}

% nguồn: Ok_TracNghiemLop11_Lan1
\begin{ex}%[1D5N1-4]%[ma:467a6ecd]
	Khảo sát thời  gian tập thể dục của một số học sinh khối $11$ thu được mẫu số liệu ghép nhóm sau:
	\begin{center}
		\begin{tabular}{|c|c|c|c|c|}
			\hline 
			 Thời gian (Phút) &$\left[0;20\right)$ & $\left[20;40\right)$ &$\left[40;60\right)$ &$\left[60;80\right)$ \\
			\hline 
			Số học sinh & $5$ &$9$ &$12$ &$10$ \\ 
			 \hline
		\end{tabular}
	\end{center}
	Nhóm chứa mốt của mẫu số  liệu trên là
	\choice 
		 {$[0;20)$}
		 {$[20;40)$}
		 {\True $[40;60)$}
		 {$[60;80)$}
	 \loigiai{
		Nhóm chứa mốt là nhóm có tần số lớn  nhất.\\
		Do đó nhóm chứa mốt là nhóm $[40;60)$ có tần số  là $12$.
	 }
\end{ex}

% nguồn: Ok_TracNghiemLop11_Lan1
\begin{ex}%[1D5N1-4]%[ma:be8e43a0]
	Khảo sát thời gian tập  thể dục của một số học sinh khối 11 thu được mẫu số liệu ghép nhóm sau:
	\begin{center}
		\begin{tabular}{|l|c|c|c|c|c|}
			\hline Thời gian  (phút) & {$[0; 20)$} & {$[20; 40)$} & {$[40; 60)$} & {$[60; 80)$} & {$[80; 100)$} \\
			\hline Số học sinh  & $5$ & $9$ & $12$ & $10$ & $6$\\
			 \hline
		\end{tabular}
	\end{center}
	Nhóm chứa mốt của mẫu số liệu trên  là
	 \choice
		{\True  $[40; 60)$}
		{$[20; 40)$} 
		{$[60;  80)$}
		{$[80; 100)$} 
	\loigiai{ 
		Mốt $M_0$ chứa trong nhóm $[40; 60)$. 
	 }
\end{ex}

% nguồn: Ok_TracNghiemLop11_Lan1
\begin{ex}%[1D5N1-4]%[ma:34e8b325]
	Bảng dữ liệu ghép nhóm dưới đây ghi lại số khách hàng nữ tương ứng từng khoảng độ tuổi đã vào  siêu thị X trong giờ đầu mở cửa hôm 9/9.
	\begin{center}
		\setlength\extrarowheight{2pt} 
		\begin{tabular}{|c|c|c|c|c|c|}
			\hline Khoảng tuổi &  $[20; 30)$ & $[30; 40)$ & $[40; 50)$ & $[50; 60)$ & $[60; 70)$ \\
			\hline Số khách hàng nữ & $5$ &$9$ & $6$ & $4$ &  $2$ \\
			 \hline
		\end{tabular}
	\end{center}
	Xác định nhóm chứa mốt của bảng dữ liệu ghép  nhóm trên.
	\choice 
		{\True $[30;40)$} 
		 {$[40;50)$}
		 {$[20;30)$}
		 {$[60;70)$}
	\loigiai{ 
		Nhóm có tần số lớn nhất trong bảng dữ liệu là $[30;40)$ và đó là nhóm chứa mốt duy  nhất.
	} 
\end{ex}

% nguồn: Ok_TracNghiemLop11_Lan1
\begin{ex}%[1D5N1-4]%[ma:cfb327f3]
	Cho mẫu số liệu về chiều cao $(\mathrm{cm})$ của các học sinh nữ trong khối 11  của một trường như sau:
	\begin{center}
		\begin{tabular}{|l|c|c|c|c|c|c|}
			\hline Chiều cao & {$[145; 150)$} & {$[150; 155)$} & {$[155;  160)$} & {$[160; 165)$} & {$[165; 170)$} & {$[170; 175)$} \\
			\hline Số học  sinh & $20$ & $45$ & $34$ & $27$ & $15$ & $4$ \\
			 \hline
		\end{tabular}
	\end{center}
	Số học sinh nữ cao từ $150  \mathrm{~cm}$ đến $155 \mathrm{~cm}$ là
	\choice 
		{$20$} 
		{$65$} 
		 {$34$}
		{\True  $45$}
	 \loigiai{
		Theo bảng số liệu trên, số học sinh nữ cao từ $150 \mathrm{~cm}$ đến $155 \mathrm{~cm}$ là $45$  học sinh.
	} 
\end{ex}

% nguồn: Ok_TracNghiemLop11_Lan1
\begin{ex}%[1D5N1-4]%[ma:80147728]
	Cho mẫu số liệu về thời gian (phút) đi từ nhà đến trường của các học sinh trong một lớp 11 của một trường như  sau:
	\begin{center}
		\begin{tabular}{|l|c|c|c|c|c|c|}
			\hline Thời gian &  {$[0; 5)$} & {$[5; 10)$} & {$[10; 15)$} & {$[15; 20)$} & {$[20; 25)$} & {$[25; 30)$} \\
			\hline Số học sinh & $7$ & $12$  & $7$ & $5$ & $3$ & $2$ \\
			\hline 
		\end{tabular}
	\end{center}
	Có bao  nhiêu học sinh có thời gian đi từ nhà đến trường là $15$ phút đến $20$ phút?
	 \choice
		{$20$} 
		{$15$} 
		 {\True $5$}
		 {$7$}
	\loigiai{ 
		Theo bảng số liệu,  có $5$ học sinh có thời gian đi từ nhà đến trường là $15$ phút đến $20$ phút.
	 }
\end{ex}

% nguồn: Ok_TracNghiemLop11_Lan1
\begin{ex}%[1D5N1-4]%[ma:790af0a9]
	Khảo sát chiều cao của một số học sinh khối $11$ thu  được mẫu số liệu ghép nhóm sau:
	\begin{center}
		\begin{tabular}{|c|c|c|c|c|c|}
			\hline Khoảng chiều cao  (cm) & {$[145; 150)$} & {$[150; 155)$} & {$[155; 160)$} & {$[160; 165)$} & {$[165; 170)$} \\
			\hline Số học sinh & $7$ & $14$ & $10$ & $10 $& $9$ \\ 
			\hline 
		\end{tabular}
	\end{center}
	 Tính mốt của mẫu số liệu ghép nhóm này (làm tròn kết quả đến hàng phần trăm)
	 \choice
		{$160$}{$152{,}25$}{\True $153{,}18$}{$170$} 
	 \loigiai{Tần số lớn nhất là 14 nên nhóm chứa mốt là nhóm $[150; 155)$.\\
		Ta có $j=2; a_2=150;  m_2=14; m_1=7; m_3=10; h=5$.\\
		Do đó $M_0=150+\dfrac{14-7}{(14-7)+(14-10)} \cdot 5 \approx 153{,}18$. 
	} 
\end{ex}

% nguồn: Ok_TracNghiemLop11_Lan1
\begin{ex}%[1D5N1-4]%[ma:868a2036]
	Gọi $i$ là nhóm có tần số lớn nhất. Gọi $u,\; g,\; n_i$ lần lượt là đầu mút trái, độ dài và tần số của nhóm $i$; $\; n_{i-1}, \; \; n_{i+1}$ lần lượt là tần số của nhóm $i-1$, nhóm $i+1$. Gọi  $M_0$ là Mốt của mẫu số liệu ghép nhóm, mệnh đề nào sau đây đúng?
	 \choice
		{\True  $M_0=u+\left(\dfrac{n_i-n_{i-1}}{2n_i-n_{i-1}-n_{i+1}} \right)\cdot g$}
		{$M_0=u+\left(\dfrac{n_i.n_{i-1}}{2n_i-n_{i+1}-n_{i-1}}  \right)\cdot g$}
		{$M_0=u.\left(\dfrac{n_i-n_{i-1}}{2n_i-n_{i-1}-n_{i+1}}  \right)-g$}
		{$M_0=u+\left(\dfrac{n_{i+1}-n_{i-1}}{2n_i-n_{i-1}-n_{i+1}} \right)\cdot  g$}
	\loigiai{ 
	Theo định nghĩa mốt $M_0=u+\left(\dfrac{n_i-n_{i-1}}{2n_i-n_{i-1}-n_{i+1}} \right)\cdot g$.} 
\end{ex}

% nguồn: Ok_TracNghiemLop11_Lan1
\begin{ex}%[1D5N1-4]%[ma:8452e2e9]
	Khảo sát thời gian chạy bộ trong một ngày của một số học sinh khối $11$ thu được mẫu  số liệu ghép nhóm sau
	\begin{center}
		\begin{tabular}{|l|c|c|c|c|c|}
			\hline Thời gian (phút) & $[0;  20)$ & $[20; 40)$& $[40; 60)$ & $[60; 80)$ & $[80; 100)$ \\
			\hline Số  học sinh & $5$ & $9$ & $12$ & $10$ & $6$ \\
			\hline 
		\end{tabular}
	\end{center}
	Nhóm chứa  mốt của mẫu số liệu này là
	\choice 
		 {$\left[20;40\right)$}
		 {\True $\left[40;60\right)$}
		{$\left[60;80\right)$} 
		 {$\left[80;100\right)$}
	\loigiai{ 
		Nhóm chứa mốt của mẫu số liệu này  là $\left[40;60\right)$.
	 }
\end{ex}

% nguồn: Ok_TracNghiemLop11_Lan1
\begin{ex}%[1D5N1-4]%[ma:95c86e6c]
	Điểm kiểm tra $15$  phút của $36$ học sinh lớp 11A được cho bởi bảng tần số ghép nhóm sau
	\begin{center}
		\begin{tabular}{|c|c|}
			 \hline
			Nhóm điểm &  Tần số \\
			 \hline
			 $\left[1;3\right)$ & $3$ \\
			 \hline
			$\left[3;5\right)$ & $2$ \\ 
			 \hline
			$\left[5;7\right)$ & $10$ \\ 
			 \hline
			$\left[7;9\right)$ & $14$  \\
			\hline 
			$\left[9;11\right)$ & $7$  \\
			 \hline
			&  $n=36$ \\
			 \hline
		\end{tabular}
	\end{center}
	Mốt của bảng ghép  lớp trên là giá trị nào sau?
	 \choice
		{\True $7{,}73$} 
		{$6{,}12$} 
		 {$5{,}09$}
		{$7{,}03$} 
	\loigiai{ 
		Áp dụng công thức tính  mốt $M_0=u+\left(\dfrac{n_i-n_{i-1}}{2n_i-n_{i-1}-n_{i+1}} \right)\cdot g$.\\
		Với $i=4$,  $u=7$, $g=2$, $n_4=14$, $n_3=10$, $n_5=7$.\\
		Ta có 
		\[M_0=7+\left(\dfrac{14-10}{2\cdot 14-10-7} \right)\cdot 2=7{,}73.\] 
	} 
\end{ex}

% nguồn: Ok_TracNghiemLop11_Lan1
\begin{ex}%[1D5N1-4]%[ma:bdb643de]
	Khảo sát thời gian (đơn vị: phút) tập  thể dục của một số học sinh khối $11$ thu được mẫu số liệu ghép nhóm sau:
	\begin{longtable}{|c|c|c|c|c|}
		\hline 
		 Thời gian (phút)&$[0;20)$&$[20;40)$&$[40;60)$&$[60;80)$\\
		 \hline
		Số học  sinh&$5$&$9$&$12$&$10$\\
		\hline 
	\end{longtable}
	Nhóm chứa mốt của mẫu số  liệu trên là
	\choice 
		 {$[0;20)$}
		{$[20;40)$} 
		{\True  $[40;60)$}
		{$[60;80)$} 
	 \loigiai{
		Nhóm chứa mốt là  nhóm có tần số lớn nhất.\\
		Do đó nhóm chứa mốt là  nhóm $[40;60)$ có tần số là $12$.
	} 
\end{ex}

% nguồn: Ok_TracNghiemLop11_Lan1
\begin{ex}%[1D5N1-4]%[ma:5a56cf15]
	Khảo sát thời gian chạy bộ trong một ngày của một số học sinh khối $11$ thu được mẫu số liệu ghép nhóm sau 
	\begin{center}
		{\tabcolsep =  2mm
		\begin{tabular}{|c|c|c|c|c|c|}
			\hline 
			Thời gian (phút) & $[0; 20)$ & $[20; 40)$ & $[40; 60)$ & $[60; 80)$ & $[80; 100)$  \\
			 \hline
			Số học  sinh & $5$ & $9$ & $12$ & $10$ & $6$\\
			\hline 
		\end{tabular}
		 }
	\end{center}
	Mẫu số  liệu ghép nhóm này có mốt là
	 \choice
		 {$59$}
		{$40$} 
		 {\True $52$}
		 {$53$}
	 \loigiai{
		Mốt  $\mathrm{M}_0$ chứa trong nhóm $[40;60)$.\\
		Do đó, $u_m=40$; $u_{m+1}=60$  $\Rightarrow u_{m+1} - u_m = 60-40=20$.\\
		$n_{m-1}=9$; $n_m=12$;  $n_{m+1}=10$.\\
		 $\mathrm{M}_0 = 40 + \dfrac{12-9}{(12-9)+(12-10)} (60-40) =52$.
	 }
\end{ex}

% nguồn: Ok_Toan11_TracNghiem_Lan2
\begin{ex}%[1D5H1-4]%[ma:b3d3d497]
	Anh Văn ghi lại cự li $30$ lần ném lao của mình ở bảng sau  (đơn vị: mét) rồi tổng hợp lại kết quả ném của anh Văn vào bảng tần số ghép nhóm theo mẫu sau:
	\begin{center}
		\begin{tabular}{|c|c|c|c|c|c|}
			 \hline
			Cự li (m) & $\left[69{,}2;70\right)$ &  $\left[70;70{,}8\right)$ & $\left[70{,}8;71{,}6\right)$ & $\left[71{,}6;72{,}4\right)$ & $\left[72{,}4;73{,}2\right)$ \\
			 \hline
			Sồ lần & $4$ & $2$ & $9$ & $10$ &  $5$ \\
			\hline 
		\end{tabular}
	\end{center}
	 Khả năng anh Văn ném được khoảng bao nhiêu mét là cao nhất?
	\choice 
		{$47{,}7$} 
		{$65{,}6$} 
		{\True $71{,}5$} 
		{$49{,}9$} 
	 \loigiai{
		\begin{itemize}
			\item Nhóm chứa  mốt của mẫu số liệu trên là $\left[71{,}6;72{,}4\right)$.
			\item Ta có:  $u_m=71{,}6$; $n_{m-1}=9$; $n_{m+1}=5$; $u_{m+1}-u_m=72{,}4-71{,}6=0{,}8$.
			\item Mốt của mẫu số liệu ghép nhóm là \[M_0=71{,}6+\dfrac{10-9}{(10-9)+(10-5)} \cdot  0{,}8=71{,}7\left(\mathrm{m}\right).\]
		\end{itemize}
		Vậy khả năng anh Văn ném được $71{,}7$ m là  cao nhất.
	 }
\end{ex}

% nguồn: Ok_Toan11_TracNghiem_Lan2
\begin{ex}%[1D5H1-4]%[ma:1920acb9]
	Kết quả điểm thi của $32$ học sinh trong kì thi Tiếng Anh được biểu diễn  ở biểu đồ dưới đây
	\begin{center}
		\begin{tikzpicture}[xscale=2,font=\footnotesize,line join=round,line cap=round,>=stealth]
				\draw[->](0,-0.2)--(0,7) node[left]{\textbf{Tỉ  lệ}};
				\draw[->](0,0)--(7.5,0) node[below]{\textbf{điểm}}; 
				\foreach \i in{0,...,6}  \pgfmathsetmacro{\gti}{int (5*(\i))}
				\draw [dotted](0,\i) circle (1pt) node[left]{$\gti\%$}  -- (5.5,\i);
				\foreach  \i/\a/\b in {1/40/50,2/50/60,3/60/70,4/70/80,5/80/90}\draw (\i+.25,0) node [below] {$[\a; \b)$};
				 \foreach \i/\a/\b in {6/90/100}\draw (\i+0.25,0) node [below] {$[\a; \b]$};
				\foreach \i/\ph  in {1/13,2/19,3/31,4/19,5/13,6/6}
					{ 
				\draw[fill=orange](\i,0) rectangle (\i+0.5,0.2*\ph) node [above left=0.1]{$\ph\%$}; 
				 }
			\end{tikzpicture}
	\end{center}
	Hãy ước lượng mốt của mẫu số liệu trên? 
	 \choice
		{$55$} 
		 {\True $65$}
		 {$72$}
		{$68$} 
	\loigiai{ 
		Ta có bảng tần  số ghép nhóm sau
		\begin{center}
			\begin{tabular}{|c|c|c|c|c|c|c|c|}
				 \hline
				 Lớp điểm& $[40;\,50)$ & $[50;\,60)$ & $[60;\,70)$ & $[70;\,80]$ & $[80;\,90]$ & $[90;\,100]$\\
				\hline 
				 Tần số & $4$ & $6$ & $10$ & $6$ & $4$ &$2$\\
				 \hline
			\end{tabular}
		\end{center}
		Nhóm chứa  mốt của mẫu số liệu là $\left[60;\,70 \right)$.\\
		Do đó $u_3=60$; $u_4=70\Rightarrow u_4-u_3=70-60=10$,\\ 
		và $n_2=6$; $n_3=10$;  $n_4=6$.\\
		Mốt của mẫu  số liệu là $60+\dfrac{10-6}{10-6+10-6} \cdot 10=65$.
	 }
\end{ex}

% nguồn: Ok_Toan11_TracNghiem_Lan2
\begin{ex}%[1D5H1-4]%[ma:8dd4f176]
	Điểm kiểm tra giữa kỳ $I$ của một  lớp được cô giáo chủ nhiệm ghi lại theo bảng sau
	\begin{center}
		\begin{tabular}{|c|c|c|c|c|c|c|c|}
			\hline 
			Điểm& $[4{,}5;5{,}5)$ & $[5{,}5;6{,}5)$ & $[6{,}5;7{,}5)$ & $[7{,}5;8{,}5)$  & $[8{,}5;9{,}5)$\\
			 \hline
			Số học sinh & $3$ & $7$ & $9$ & $15$ &  $6$ \\
			 \hline
		\end{tabular}
	\end{center}
	Mốt của số  liệu ghép nhóm trên là
	\choice 
		 {$8{,}4$}
		 {$8{,}9$}
		{\True  $7{,}9$}
		{$11{,}6$} 
	 \loigiai{
		Nhóm chứa mốt của  mẫu số liệu là $\left[7{,}5;8{,}5 \right)$.\\
		Do đó $u_4=7{,}5$;  $u_5=8{,}5\Rightarrow u_5-u_4=8{,}5-7{,}5=1$,\\
		và  $n_3=9$; $n_4=15$; $n_5=6$.\\
		Mốt của số liệu ghép nhóm  là $7{,}5+\dfrac{15-9}{15-9+15-6} \cdot 1=7{,}9$.
	 }
\end{ex}

% nguồn: Ok_Toan11_TracNghiem_Lan2
\begin{ex}%[1D5H1-4]%[ma:cacaeb6d]
	Kiểm tra cân nặng $25$ quả dưa của hai nhà vườn $A$ và $B$ cho kết quả như  bảng sau:
	\begin{center}
		\begin{tabular}{|c|c|c|c|c|c|c|c|}
			\hline 
			Cân nặng (kg)& $[1{,}5;2{,}5)$ & $[2{,}5;3{,}5)$  & $[3{,}5;4{,}5)$ & $[4{,}5;5{,}5)$ & $[5{,}5;6{,}5)$ \\
			 \hline
			Số quả  của vườn $A$ & $2$ & $6$ & $4$ & $10$ & $3$\\
			 \hline
			Số quả của vườn $B$ & $3$  & $7$ & $8$ & $5$ & $2$\\
			\hline 
		\end{tabular}
	\end{center}
	Nếu lấy một quả dưa bất kỳ của vườn $A$ thì khả năng cao nhất là lấy được quả dưa nặng bao nhiêu  kg?
	\choice 
		 {$3{,}84$ kg}
		 {$4{,}84$ kg}
		{$4{,}9$  kg}
		{\True $4{,}96$ kg} 
	 \loigiai{
		Nhóm chứa mốt của mẫu  số liệu là $\left[4{,}5;5{,}5 \right)$.\\
		 Do đó $u_4=4{,}5$; $u_5=5{,}5\Rightarrow u_5-u_4=5{,}5-4{,}5=1$,\\
		 và $n_3=4$; $n_4=10$; $n_5=5$.\\
		Mốt của cân nặng của các quả dưa vườn  $A$ là\\
		$M_0=4{,}5+\dfrac{10-4}{10-4+10-3}  \cdot 1=4{,}96$.\\
		 Dựa vào kết quả trên, khả năng lấy được quả dưa nặng $4{,}96$ kg là cao nhất.
	} 
\end{ex}

% nguồn: Ok_Toan11_TracNghiem_Lan2
\begin{ex}%[1D5H1-4]%[ma:318c3d52]
	Cho mẫu số liệu ghép nhóm về chiều  cao của $25$ cây dừa giống như sau
	\begin{center}
		\begin{tabular}{|c|c|c|c|c|c|c|c|}
			 \hline
			Chiều cao (cm)& $[0;10)$ & $[10;20)$ & $[20;30)$ & $[30;40)$ & $[40;50)$  \\
			\hline 
			Số cây & $4$ &  $6$ & $7$ & $5$ & $3$\\
			\hline 
		\end{tabular}
	\end{center}
	Mốt của mẫu số  liệu ghép nhóm này là
	 \choice
		{\True  $M_0=\dfrac{70}{3}$}
		 {$M_0=\dfrac{50}{3}$}
		 {$M_0=\dfrac{70}{3}$}
		 {$M_0=\dfrac{80}{3}$}
	\loigiai{ 
		Tấn  số lớn nhất là $7$ nên nhóm chứa mốt là $\left[20;30 \right)$.\\
		Do đó $u_3=20$;  $u_4=30\Rightarrow u_4-u_3=30-20=10$,\\
		và $n_2=6$;  $n_3=7$; $n_4=5$.\\
		Vậy 
		$M_0=20+\dfrac{7-6}{\left( 7-6 \right)+\left( 7-5  \right)} \cdot 10=\dfrac{70}{3}$.
	 }
\end{ex}

% nguồn: Ok_Toan11_TracNghiem_Lan2
\begin{ex}%[1D5H1-4]%[ma:08aeebfc]
	Điểm thi môn Toán (thang điểm $100$, điểm được làm tròn đến hàng phần  chục) của $60$ thí sinh được cho trong bảng sau
	\begin{center}
		\begin{tabular}{|l|c|c|c|c|c|c|c|c|c|c|c|}
			\hline Điểm&  {$[0;10)$} & {$10;20)$} & {$[20;30) $} & {$[30;40)$} & {$[40;50)$} & {$[50;60)$}& {$[60;70)$}& {$[70;80)$}& {$[80;90)$}& {$[90;100)$}\\
			\hline  Số thí sinh & $1$ & $2$ & $3$ & $6$& $15$& $12$ & $10$ & $6$& $3$& $1$\\
			 \hline
		\end{tabular}
	\end{center}
	mốt  của mẫu số liệu trên bằng
	 \choice
		{ $45{,}5$} 
		{  $46{,}5$}
		 { $45$}
		 { \True $47{,}5$}
	\loigiai{ 
		Nhóm chứa mốt của mẫu số liệu là $[40;50)$.\\ 
		 Do đó $u_m=40,n_{m-1}=6,n_{m+1}=12,u_{m+1}-u_m=50-40=10$.\\
	mốt của mẫu số liệu  ghép nhóm là $M_o=40+\dfrac{15-6}{(15-6)+(15-12)}\cdot 10=47{,}5$.}
\end{ex}

% nguồn: Ok_Toan11_TracNghiem_Lan2
\begin{ex}%[1D5H1-4]%[ma:654017ff]
	Cho mẫu số liệu ghép nhóm về thống kê nhiệt độ tại một địa điểm trong $30$ ngày,  ta có bảng số liệu sau:
	\begin{center}
		\begin{tabular}{|c|c|c|c|c|c|c|c|}
			\hline 
			Nhiệt độ (${}^\circ C$)& $[18;21)$ & $[21;24)$ & $[24;27)$ &  $[27;30)$ \\
			\hline 
			Số  ngày & $6$ & $12$ & $9$ & $3$ \\
			 \hline
		\end{tabular}
	\end{center}
	Mốt của mẫu số liệu ghép nhóm trên  là
	 \choice
		 {$24$}
		 {$25$}
		{$22$} 
		{\True $23$} 
	\loigiai{ 
		Nhóm $2$ ứng với nửa  khoảng $\left[21; 24 \right)$ là nhóm có tần số lớn nhất với $u_2=21$; $u_3=24$\\
		$\Rightarrow u_3-u_2=24-21=3$  và $n_1=6$; $n_2=12$; $n_3=9$.\\
		Ta có mốt  của mẫu số liệu là
		$M_0=21+\left( \dfrac{12-6}{2  \cdot 12-6-9} \right) \cdot 3=23$.
	} 
\end{ex}

% nguồn: Ok_Toan11_TracNghiem_Lan2
\begin{ex}%[1D5H1-4]%[ma:0a905c1c]
	Một thư viện thống kê số lượng sách được mượn mỗi ngày trong $3$ tháng ở bảng  sau
	\begin{center}
		\begin{tabular}{|l|c|c|c|c|c|c|c|c|}
			\hline Số sách& {$[16;20)$}  & {$ [21;25)$} & {$[26;30) $} & {$[31;35)$} & {$[36;40)$} & {$[41;45)$} & {$[46;50)$}\\
			\hline Số  ngày& $3$ & $6$ & $15$ & $27$& $22$& $14$& $5$ \\
			 \hline
		\end{tabular}
	\end{center}
	Trong $3$ tháng trên số sách được mượn mỗi ngày cao nhất gần  bằng bao nhiêu (kết quả làm tròn đến hàng đơn vị).
	\choice 
		{  $32$}
		{  \True $34$}
		{  $33$}
		{  $35$}
	\loigiai{ 
		Do số sách là số nguyên nên ta hiệu chỉnh lại như sau 
		\begin{center}
			\begin{tabular}{|l|c|c|c|c|c|c|c|c|}
				\hline Số sách& {$[15{,}5;20{,}5)$} & {$ [20{,}5;25{,}5)$} &  {$[25{,}5;30{,}5) $} & {$[30{,}5;35{,}5)$} & {$[35{,}5;40{,}5)$} & {$[40{,}5;45{,}5)$} & {$[45{,}5;50{,}5)$}\\
				\hline Số ngày& $3$ & $6$ & $15$ & $27$& $22$& $14$&  $5$ \\
				\hline 
			\end{tabular}
		\end{center}
		Ta có nhóm chứa  mốt của mẫu số liệu trên là nhóm $[30,5;35,5)$.\\
		Do đó $u_m=30{,}5$, $n_{m-1}=15$,  $n_m=27$, $n_{m+1}=22$, $u_{m+1}-u_m=35{,}5-30{,}5=5$.\\
	mốt của mẫu  số liệu ghép nhóm là $$M_o=30{,}5+\dfrac{27-15}{(27-15)+(27-22)}\cdot 5=\dfrac{1157}{34}\approx 34.$$}
\end{ex}

% nguồn: Ok_Toan11_TracNghiem_Lan2
\begin{ex}%[1D5H1-4]%[ma:a74ac505]
	Khảo sát thời gian tập thể dục của một số học sinh khối $11$ thu được mẫu số liệu ghép nhóm sau 
	\begin{center}
		\begin{tabular}{|c|c|c|c|c|c|}
			\hline 
			Thời  gian (phút) & $[0;20)$ & $[20;40)$ & $[40;60)$ & $[60;80)$ & $[80;100)$ \\
			\hline 
			Số học sinh & $5$  & $9$ & $12$ & $10$ & $6$ \\
			 \hline
		\end{tabular}
	\end{center}
	Mốt của mẫu số liệu  trên là
	 \choice
		{$42$} 
		 {\True $52$}
		 {$53$}
		 {$54$}
	\loigiai{ 
		Mốt  $M_0$ chứa trong nhóm $[40;60)$.\\
		 Do đó $u_3=40$; $u_4=60\Rightarrow u_4-3=60-40=20$.\\
		$n_2=9$; $n_3=12$; $n_4=10$.\\ 
		Vậy  $M_0=40+\dfrac{12-9}{(12-9)+(12-10)}(60-20)=52$.
	} 
\end{ex}

% nguồn: Ok_Toan11_TracNghiem_Lan2
\begin{ex}%[1D5H1-4]%[ma:b108fd63]
	Khảo sát thời gian tự học bài ở  nhà của một số em học sinh lớp $11$ thu được mẫu ghép nhóm số lượng như sau:
	\begin{center}
		\begin{tabular}{|c|c|c|c|c|c|c|c|}
			\hline 
			Thời gian(phút)& $[0;30)$ & $[30;60)$ & $[60;90)$  & $[90;120)$ & $[120;150)$ \\
			\hline 
			Số học sinh & $9$ & $10$ & $9$ & $15$ & $7$\\ 
			 \hline
		\end{tabular}
	\end{center}
	Mốt của số liệu ghép nhóm trên (kết quả được làm tròn đến hàng  phần trăm) là
	\choice 
		{\True  $102{,}86$}
		 {$102{,}85$}
		 {$102{,}8$}
		 {$102{,}9$}
	\loigiai{ 
		Nhóm chứa mốt  của mẫu số liệu là $\left[90;120 \right)$.\\
		Do đó $u_4=90$; $u_5=120\Rightarrow u_5-u_4=120-90=30$,\\ 
		và $n_3=9$;  $n_4=15$; $n_5=7$.\\
		 Vậy $M_0=90+\dfrac{15-9}{15-9+15-7} \cdot 30\approx 102{,}86$.
	} 
\end{ex}

% nguồn: Ok_Toan11_TracNghiem_Lan2
\begin{ex}%[1D5H1-4]%[ma:3f905271]
	Bảng số liệu ghép nhóm sau cho biết chiều cao của $50$ học  sinh $11A$ như sau
	\begin{center}
		\begin{tabular}{|c|c|c|c|c|c|c|c|}
			 \hline
			Khoảng chiều cao & $[145;150)$ &  $[150;155)$ & $[155;160)$ & $[160;165)$ & $[165;170)$ \\
			\hline 
			Số học  sinh & $7$ & $14$ & $10$ & $10$ & $9$\\
			 \hline
		\end{tabular}
	\end{center}
	Tính mốt của  số liệu ghép nhóm này.
	\choice 
		{\True $153{,}18$} 
		{$152{,}18$} 
		{$154{,}01$} 
		 {$156{,}72$}
	\loigiai{ 
		Tần số lớn nhất là $14$ nên nhóm chứa mốt là  nhóm $\left[150;155 \right)$.\\
		Do đó $u_2=150$; $u_3=155$ $\Rightarrow u_3-u_2=155-150=5$ và  $n_1=7$; $n_2=14$; $n_3=10$.\\
		Do đó $M_0=150+\dfrac{14-7}{\left( 14-7 \right)+\left( 14-10 \right)}\cdot 5\approx  153{,}18$.
	} 
\end{ex}

% nguồn: Ok_Toan11_TracNghiem_Lan2
\begin{ex}%[1D5H1-4]%[ma:69da2f2d]
	Thời gian (phút) đề học sinh hoàn thành một  câu hỏi thi được cho như sau
	\begin{center}
		\begin{tabular}{|l|c|c|c|c|c|c|c|}
			\hline Thời gian & {$[0{,}5;10{,}5)$} & {$ [10{,}5;20{,}5)$} & {$[20{,}5;30{,}5) $} & {$[30{,}5;40{,}5)$}  & {$[40{,}5;50{,}5)$}\\
			\hline Số học sinh & $2$ & $10$ & $6$  & $4$& $3$\\
			\hline 
		\end{tabular}
	\end{center}
	Tìm mốt của mẫu số liệu  ghép nhóm này.
	\choice 
		{  $17{,}42$}
		{  $14{,}56$}
		 { \True $17{,}16$}
		 { $12{,}67$}
	 \loigiai{
		Tần số lớn nhất là $10$ nên nhóm chứa mốt  là nhóm $[10{,}5;20{,}5]$.\\
		Ta có $u_m=10{,}5,u_{m+1}=20{,}5,n_m=10;n_{m+1}=6,n_{m-1}=2,u_{m+1}-u_m=10.$\\ 
	Do đó  $M_o=10{,}5+\dfrac{10-2}{(10-2)+(10-6)}\cdot 10=17{,}16$.}
\end{ex}

% nguồn: Ok_Toan11_TracNghiem_Lan2
\begin{ex}%[1D5H1-4]%[ma:0ddc7eef]
	 Doanh thu bán hàng trong $20$ ngày được lựa chọn ngẫu nhiên của một của hàng được ghi lại ở bảng sau (đơn vị: triệu đồng):
	\begin{center}
		\begin{tabular}{|c|c|c|c|c|c|}
			\hline 
			Doanh thu  & $[5;7)$ & $[7;9)$ & $[9;11)$ & $[11;13)$ & $[13;15)$ \\
			\hline 
			 Số ngày & $2$ & $7$ & $7$ & $3$ & $1$ \\
			\hline 
		\end{tabular}
	\end{center}
	Mốt của mẫu số liệu trên thuộc khoảng nào trong  các khoảng dưới đây?
	\choice 
		{$\left[7;9  \right)$}
		{\True $\left[9;11  \right)$}
		{$\left[11;13  \right)$}
		{$\left[13;15  \right)$}
	\loigiai{ 
		Có $2$ nhóm chứa mốt của mẫu số liệu trên  đó là $\left[7;9 \right)$ và $\left[9;11 \right)$.\\
		Xét nhóm $\left[7;9 \right)$  ta có $u_2=7$; $u_{3}=9\Rightarrow u_{3}-{u_2}=9-7=2$,\\
		và  $n_{1}=2$; $n_2=7$; $n_{3}=7$.\\
		Suy ra  $M_0=7+\dfrac{7-2}{(7-2)+(7-7)}(9-7)=9$.\\
		Xét  nhóm $\left[9;11 \right)$ ta có $u_3=9$; $u_4=11\Rightarrow u_4-u_3=11-9=2$,\\
		và  $n_2=7$; $n_3=7$; $n_4=3$.\\ $M'_0=9+\dfrac{7-7}{(7-7)+(7-3)}(11-9)=9$.\\
		 Vậy mốt của mẫu số liệu là $9$.
	} 
\end{ex}

% nguồn: Ok_Toan11_TracNghiem_Lan2
\begin{ex}%[1D5H1-4]%[ma:ceb3a7e8]
	Tuổi thọ (năm) của $50$ bình ắc quy ô tô  được cho như sau:
	\begin{center}
		\begin{tabular}{|c|c|c|c|c|c|c|}
			 \hline
			Tuổi thọ  (năm) & $[2;2{,}5)$ & $[2{,}5;3)$ & $[3;3{,}5)$ & $[3{,}5;4)$ & $[4;4{,}5)$ & $[4{,}5;5)$ \\
			 \hline
			Tần số & $4$  & $9$ & $14$ & $11$ & $7$ & $5$\\
			\hline 
		\end{tabular}
	\end{center}
	 Xác định mốt?
	 \choice
		{\True $3{,}3125$} 
		 {$3{,}2645$}
		{$3{,}4534$} 
		{$2{,}9345$} 
	\loigiai{ 
		Nhóm chứa mốt là $\left[3;3{,}5 \right)$.\\ 
		 Do đó $u_3=3$; $u_4=3{,}5\Rightarrow u_4-3=3{,}5-3=0{,}5$, \\
		và $n_2=9$; $n_3=14$;  $n_4=11$.\\
		$M_0=3+\dfrac{14-9}{14-9+14-11} \cdot 0{,}5=3{,}3125$ 
	} 
\end{ex}

% nguồn: Ok_Toan11_TracNghiem_Lan2
\begin{ex}%[1D5H1-4]%[ma:4a354372]
	Thời gian (phút) truy cập Internet mỗi buổi tối  của một số học sinh được cho trong bảng sau
	\begin{center}
		\begin{tabular}{|l|c|c|c|c|c|c|c|}
			\hline Thời gian & {$[9{,}5;12{,}5)$} & {$ [12{,}5;15{,}5)$} & {$[15{,}5;18{,}5) $} & {$[18{,}5;21{,}5)$}  & {$[21{,}5;24{,}5)$}\\
			\hline Số học sinh & $3$ & $12$ & $15$ & $24$& $2$\\ 
			\hline 
		\end{tabular}
	\end{center}
	Tìm mốt  của mẫu số liệu ghép nhóm này (kết quả làm tròn đến hàng phần trăm).
	\choice 
		{ $18{,}37$} 
		 { $19{,}37$}
		{ \True $19{,}73$} 
		 { $18{,}73$}
	\loigiai{ 
		Tần số lớn nhất là $24$  nên nhóm chứa mốt là nhóm $[18{,}5;21{,}5]$.\\
		 Ta có: $u_m=18{,}5,u_{m+1}=21{,}5,n_m=24,n_{m+1}=15,n_{m-1}=2,u_{m+1}-u_m=3.$\\
	mốt của mẫu số  liệu là $M_o=18,5+\dfrac{24-15}{(24-15)+(24-2)}\cdot 3=18{,}5+\dfrac{27}{31}\approx 19{,}37$.}
\end{ex}

% nguồn: Ok_Toan11_TracNghiem_Lan2
\begin{ex}%[1D5H1-4]%[ma:575c4246]
	Số  sản phẩm bán được của mặt hàng $X$ tại một cửa hàng trong $30$ ngày được thống kê trong bảng
	\begin{center}
		\begin{tabular}{|c|c|c|c|c|c|c|}
			 \hline
			Số sản phẩm & $[3;5]$ & $[6;8]$ & $[9;11]$ & $[12;14]$ & $[15;17]$  \\
			\hline 
			Số học sinh & $3$ & $5$ &  $7$ & $10$ & $5$\\
			\hline 
		\end{tabular}
	\end{center}
	Cửa hàng đó bán được cao nhất  bao nhiêu sản phẩm mỗi ngày?
	\choice 
		 {\True $12$}
		{$13$} 
		 {$14$}
		 {$17$}
	 \loigiai{
		Nhóm chứa mốt  của mẫu số liệu là $\left[12;\,14 \right]$.\\
		Do  đó $u_4=12$; $u_5=14\Rightarrow u_5-u_4=14-12=2$,\\
		và $n_3=7$; $n_4=10$;  $n_5=5$.\\
		Mốt của mẫu số liệu là $12+\dfrac{10-7}{10-7+10-5} \cdot  2$$=12{,}75$.\\
		Vậy  cửa hàng đó bán được cao nhất $12$ sản phẩm mỗi ngày.
	} 
\end{ex}

% nguồn: Ok_TracNghiemLop11_Lan1
\begin{ex}%[1D5H1-4]%[ma:57b54b12]
	Khảo sát thời gian tập thể dục trong ngày của một số học sinh khối 11 thu được mẫu số liệu ghép nhóm  sau
	\begin{center}
		\begin{tabular}{|c|c|c|c|c|c|}
			\hline 
			Thời  gian (phút)&$[0;20)$&$[20;40)$&$[40;60)$&$[60;80)$&$[80;100)$\\
			 \hline
			Số học sinh  &$5$&$9$&$12$&$10$&$6$\\
			 \hline
		\end{tabular}
	\end{center}
	Nhóm chứa  mốt của mẫu số liệu này là
	 \choice
		 {$[20;40)$}
		 {\True $[40;60)$}
		 {$[60;80)$}
		 {$[80;100)$}
	\loigiai{Tần số lớn nhất là $12$ nên nhóm  chứa mốt là nhóm $[40;60)$.}
\end{ex}

% nguồn: Ok_TracNghiemLop11_Lan1
\begin{ex}%[1D5H1-4]%[ma:95460a67]
	An tìm hiểu hàm lượng chất béo (đơn vị: g) có trong $100$g mỗi loại thực phẩm. Sau khi thu thập dữ liệu về $60$ loại thực phẩm, An lập được bảng thống kê. Tìm mốt  của mẫu số liệu.
	\begin{center}
		\begin{tabular}{|c|c|c|c|c|c|c|}
			\hline Hàm lượng chất béo g & {$[2; 6)$}  & {$[6; 10)$} & {$[10; 14)$} & {$[14; 18)$} & {$[18; 22)$} & {$[22; 26)$} \\
			\hline Tần số & $2$ & $6$ & $10$ & $13$ &  $16$ & $13$ \\
			\hline 
		\end{tabular}
	\end{center}
	\choice 
		{$10$} 
		 {\True $20$}
		{$97$} 
		 {$95$}
	 \loigiai{
		Nhóm chứa mốt của mẫu  số liệu trên là nhóm $[18; 22)$.\\
		Do đó $u_m=18$, $n_{m-1}=13$, $n_m=16$,  $n_{m+1}=13$, $u_{m+1}-u_m=22-18=4$.\\
		Mốt  của mẫu số liệu ghép nhóm là
		 $$
		M_o=18+\dfrac{16-13}{(16-13)+(16-13)}\cdot 4=20. 
		 $$
	} 
\end{ex}

% nguồn: Ok_TracNghiemLop11_Lan1
\begin{ex}%[1D5H1-4]%[ma:52da808a]
	 Bảng số liệu ghép nhóm sau cho biết chiều cao (cm) của $50$ học sinh lớp $11A$
	\begin{center}
		\begin{tabular}{|l|c|c|c|c|c|}
			\hline Khoảng chiều cao (cm) & {$ [145; 150) $} & {$ [150; 155) $} & {$ [155; 160) $} & {$ [160;  165) $} & {$ [165; 170) $} \\
			\hline Số học sinh & $7$ & $14$ & $10$ & $10$ & $9$ \\ 
			\hline 
		\end{tabular}
	\end{center}
	Tính mốt  của mẫu số liệu ghép nhóm này (làm tròn đến hàng phần trăm).
	\choice 
		{\True  $153,18$}
		{$153,81$} 
		{$154,18$} 
		 {$153,28$}
	\loigiai{ 
		Tần số lớn nhất là  $14$ nên nhóm chứa mốt là nhóm $[150; 155)$.\\
		Ta có $j=2; a_{2}=150; m_{2}=14; m_{1}=7;  m_{3}=10; h=5 $.\\
		Do đó $\mathrm{M}_{0}=150+\dfrac{14-7}{(14-7)+(14-10)}  \cdot 5 \approx 153{,}18$.
	} 
\end{ex}

% nguồn: Ok_TracNghiemLop11_Lan1
\begin{ex}%[1D5H1-4]%[ma:82adf563]
	Anh Văn ghi lại cự li  $30$ lần ném lao của mình ở bảng sau (đơn vị: mét) rồi tổng hợp lại kết quả ném của anh Văn vào bảng tần số ghép nhóm theo mẫu sau
	\begin{center}
		\begin{tabular}{|l|c|c|c|c|c|}
			\hline  Cự li (m)&{$[69{,}2; 70)$}&{$[70; 70{,}8)$}&{$[70{,}8; 71{,}6)$}&{$ [71{,}6; 72{,}4)$}&{$[72{,}4; 73{,}2)$} \\
			\hline Sồ lần & $4$ &  $2$ & $9$ & $10$ & $5$ \\
			 \hline
		\end{tabular}
	\end{center}
	Khả năng anh Văn  ném được khoảng bao nhiêu mét là cao nhất?
	 \choice
		 {$47{,}7$}
		 {$65{,}6$}
		{\True $71{,}5$} 
		{$49{,}9$} 
	 \loigiai{
		Nhóm chứa mốt của mẫu  số liệu trên là $[71{,}6; 72{,}4)$.\\
		Do đó: $u_{m}=71{,}6; n_{m-1}=9; n_{m+1}=5;  u_{m+1}-u_{m}=72{,}4-71{,}6=0{,}8 $.\\
		Mốt của mẫu số liệu ghép nhóm là $\mathrm{M}_{0}=71{,}6+\dfrac{10-9}{(10-9)+(10-5)} \cdot 0{,}8=71{,}7$  (m).\\
		 Vậy khả năng anh Văn ném được $71{,}7$ m là cao nhất.
	 }
\end{ex}

% nguồn: Ok_TracNghiemLop11_Lan1
\begin{ex}%[1D5H1-4]%[ma:a626d653]
	Thống kê về thời lượng (đơn vị: phút) mỗi trận đấu bi-a trong vòng  tứ kết giải đấu European Open, người ta thu được mẫu số liệu ghép nhóm như sau:
	\begin{center}
		\begin{tabular}{|c|c|c|c|c|c|}
			\hline 
			Thời lượng & {$[9; 12)$} & {$[12; 15)$}  & {$[15; 18)$} & {$[18; 21)$} & {$[21; 24)$} \\
			 \hline
			Số trận & $3$& $12$& $14$& $24$& $4$\\ 
			 \hline
		\end{tabular}
	\end{center}
	Mốt của  mẫu số liệu ghép nhóm trên là
	\choice 
		 {$22$ phút}
		{$21$  phút}
		 {$19$ phút}
		{\True $18$  phút}
	\loigiai{ 
		Từ bảng trên, ta thấy khoảng thời lượng có số trận nhiều nhất là khoảng $[18;21)$  với $24$ trận.\\
		Vậy, mốt của mẫu số  liệu ghép nhóm trên là $18$ phút.
	} 
\end{ex}

% nguồn: Ok_TracNghiemLop11_Lan1
\begin{ex}%[1D5H1-4]%[ma:45f6f290]
	 Số khách hàng nam mua bảo hiểm ở từng độ tuổi được thống kê như sau
	\begin{center}
		{\tabcolsep  = 2mm
		\begin{tabular}{|c|c|c|c|c|c|}
			\hline 
			\textbf{Độ tuổi} & $[20; 30)$ & $[30; 40)$ & $[40; 50)$  & $[50; 60)$ & $[60; 70)$ \\
			\hline \textbf{Số khách hàng nam} & $4$ & $6$ &  $10$ & $7$ & $3$\\
			\hline 
		\end{tabular}
		} 
	\end{center}
	Hãy sử dụng dữ liệu ở trên để tư vấn cho đại lí bảo hiểm xác định khách hàng nam ở tuổi nào hay mua bảo hiểm nhất. 
	 \choice
		 {$47$}
		{\True $46$} 
		 {$48$}
		{$49$} 
	\loigiai{ 
		Nhóm chứa mốt của mẫu số liệu  khách hàng nam là $[40; 50)$.\\
		Do đó, $u_m=40$;  $u_{m+1}=50$ $\Rightarrow u_{m+1} - u_m = 50-40=10$.\\
		 $n_{m-1}=6$; $n_m=10$; $n_{m+1}=7$.\\
		$\mathrm{M}_0 = 40 + \dfrac{10-6}{(10-6)+(10-7)} (50-40) =45{,}7$.\\ 
		Dựa vào kết quả trên ta có thể  dự đoán được khách hàng nam $46$ tuổi có nhu cầu mua bảo hiểm cao nhất.
	 }
\end{ex}

% nguồn: Ok_TracNghiemLop11_Lan1
\begin{ex}%[1D5H1-4]%[ma:1bd6c68e]
	Người ta tiến hành phỏng vấn $40$ người về một mẫu áo khoác. Người điều  tra yêu cầu cho điểm mẫu áo đó theo thang điểm là $100$. Kết quả được trình bày trong bảng ghép nhóm sau
	\begin{center}
		\begin{tabular}{|c|c|c|c|c|c|c|}
			 \hline
			Nhóm  & $\left[50;60\right)$ & $\left[60;70\right)$ & $\left[70;80\right)$ & $\left[80;90\right)$ & $\left[90;100\right)$ & \\
			\hline 
			Tần số & $4$ & $5$ & $23$ & $6$ & $2$ & $N=40$  \\
			\hline 
		\end{tabular}
	\end{center}
	Mốt của mẫu số liệu ghép nhóm trên (làm tròn kết quả đến hàng đơn  vị) là
	 \choice
		 {$73$}
		{\True  $74$}
		{$75$} 
		 {$76$}
	 \loigiai{
		Nhóm chứa  mốt của mẫu số liệu trên là $[70; 80)$.\\
		Do  đó: $u_m=70 $; $ n_{m-1}=5 $; $ n_m=23$; $ n_{m+1}=6 $; $ u_{m+1}-u_m=10$.\\
		Ta có  $M_0=70+\dfrac{18}{18+27}\cdot 10\approx74$.
	 }
\end{ex}

% nguồn: Ok_TracNghiemLop11_Lan1
\begin{ex}%[1D5H1-4]%[ma:36b4fd48]
	Điểm kiểm tra giữa học kì  $I$ của lớp $11A$ được thống kê theo bảng sau:
	\begin{center}
		\begin{tabular}{|c|c|c|c|c|c|}
			\hline Số điểm & { $[0; 2)$ }  & { $[2; 4)$ } & { $[4; 6)$ } & { $[6; 8)$ } & { $[8; 10]$ } \\
			\hline Số học sinh &  $1$ & $3$ & $8$ & $18$ & $10$ \\
			\hline 
		\end{tabular}
	\end{center}
	Trong các nhóm sau  đây, nhóm nào chứa mốt của bảng số liệu ghép nhóm trên?
	\choice 
		 {$[8; 10]$}
		{$[4;  6)$}
		{\True $[6; 8)$} 
		{$[2; 4)$} 
	\loigiai{ 
		Nhóm $[6; 8)$ có tần số lớn nhất nên là  nhóm chứa mốt.
	} 
\end{ex}

% nguồn: Ok_TracNghiemLop11_Lan1
\begin{ex}%[1D5H1-4]%[ma:2119336c]
	Một công ty xây dựng khảo sát khách hàng xem họ có nhu cầu mua nhà ở mức giá  nào. Kết quả khảo sát được ghi lại ở bảng sau
	\begin{center}
		\begin{tabular}{|c|c|c|c|c|c|}
			\hline \begin{tabular}{c}
			\textbf{Mức  giá} \\
			 \textbf{(triệu đồng/$\mathrm{m}^2)$}
		\end{tabular} &{$[10; 14)$} &{$[14; 18)$} &{$[18; 22)$} &{$[22; 26)$} &{$[26; 30)$} \\
		\hline \textbf{Số khách hàng} & 54 & 78 & 120 & 45  & 12 \\
		 \hline
	\end{tabular}
\end{center}
Tìm mốt  của mẫu số liệu ghép nhóm trên (làm tròn đến hàng phần mười).
 \choice
	{$18{,}3$} 
	{\True  $19{,}4$}
	{$12{,}4$} 
	{$15{,}5$} 
 \loigiai{
	Nhóm chứa mốt của mẫu số liệu trên là nhóm $[18; 22)$.\\ Do đó $u_m=18$, $n_{m-1}=78$, $n_m=120$, $n_{m+1}=45$, $u_{m+1}-u_m=22-18=4$.\\ 
	Mốt của mẫu số liệu ghép  nhóm là
	\[M_0=18+\dfrac{120-78}{(120-78)+(120-45)} \cdot 4=\dfrac{758}{39} \approx 19{,}4. \] 
 }
\end{ex}

% nguồn: Ok_TracNghiemLop11_Lan1
\begin{ex}%[1D5H1-4]%[ma:bc75ea02]
	Người ta tiến hành phỏng vấn $40$ người về một mẫu áo sơ mi mới. Người điều tra yêu cầu cho điểm mẫu áo đó theo thang điểm là $100$. Kết  quả được trình bày trong \textit{Bảng dưới}.
	\begin{center}
		\begin{tabular}{|c|c|c|}
			\hline 
			Nhóm & Tần số & Tần số  tích lũy \\
			 \hline
			[50;60) & 4 & 4 \\ 
			\hline 
			[60;70) &  5 & 9 \\
			 \hline
			[70;80) & 23  & 32 \\
			 \hline
			[80;90)  & 6 & 38 \\
			 \hline
			[90;100) & 2 &  40 \\
			\hline 
			\multicolumn{3}{|c|}{$n=40$}  \\
			\hline 
		\end{tabular}
	\end{center}
	Mốt của mẫu số liệu ghép nhóm trên (làm tròn đến kết quả hàng đơn vị)  là
	\choice 
		{$73$} 
		 {$74$}
		{\True $75$} 
		 {$76$}
	\loigiai{ 
		Ta thấy nhóm $3$ ứng với nửa khoảng $[70;80)$ là nhóm có tần số lớn nhất với $u=70;g=10;n_3=23$.  Nhóm $2$ có tần số $n_2=5$, nhóm $4$ có tần số $n_4=6$.\\
		Áp dụng công thức, ta có mốt của mẫu số liệu  là:
	 $$ M_o=70+\left(\dfrac{23-5}{2\cdot23-5-6}\right)\cdot10 \approx 75.$$}
\end{ex}

% nguồn: Ok_TracNghiemLop11_Lan1
\begin{ex}%[1D5H1-4]%[ma:5926dcf7]
	Mốt của mẫu số liệu trên thuộc khoảng nào trong các  khoảng dưới đây?
	 \choice
		 {$[7;9)$}
		 {\True $[9;11)$}
		{$[11;13)$} 
		 {$[13;15)$}
	\loigiai{ 
		Có 2 nhóm chứa mốt của mẫu số liệu trên, do  đó
		\begin{itemize}
			\item Xét  nhóm $[7;9)$:
			 \[M_0=7+\dfrac{7-2}{(7-2)+(7-7)}\cdot (9-7)=9. \]
			\item Xét nhóm  $[9;11)$:
			\[M_0^\prime=7+\dfrac{7-7}{(7-7)+(7-3)}\cdot (11-9)=9.  \]
		\end{itemize}
		Do đó Mốt của mẫu số liệu  bằng $9$.
	 }
\end{ex}

% nguồn: Ok_TracNghiemLop11_Lan1
\begin{ex}%[1D5H1-4]%[ma:c91b2806]
	Mẫu số  liệu ghép nhóm với tần số các nhóm bằng nhau có số mốt là
	 \choice
		{\True $0$} 
		 {$1$}
		{$2$} 
		 {$3$}
	\loigiai{ 
		Khi tần số của các nhóm số liệu bằng nhau thì mẫu số liệu ghép nhóm  không có mốt.
	} 
\end{ex}

% nguồn: Ok_TracNghiemLop11_Lan1
\begin{ex}%[1D5V1-4]%[ma:88277d2e]
	Cho bảng mẫu số liệu ghép nhóm là chiều cao của  học sinh lớp $5$ tuổi như sau ($x$ nguyên dương)
	\begin{center}
		\begin{tabular}{|c|c|}
			\hline 
			Nhóm chiều cao & Tần số \\ 
			\hline 
			$\left[85;90\right)$ & $1$  \\
			 \hline
			$\left[90;95\right)$  & $x^2+5$ \\
			\hline 
			$\left[95;100\right)$ & $4x$ \\ 
			\hline 
			$\left[100;105\right)$ & $12$  \\
			 \hline
			 $\left[105;110\right)$ & $3$ \\
			\hline 
			 $\left[110;115\right)$ & $2$ \\
			\hline 
		\end{tabular}
	\end{center}
	Tìm giá trị  $x$, biết mốt của bảng ghép lớp trên phân bố $\left[90;95\right)$ là $\dfrac{283}{3}$?
	\choice 
		{\True $3$} 
		 {$4$}
		 {$5$}
		 {$6$}
	 \loigiai{
		Ta có,  theo công thức tính mốt thì
		\begin{eqnarray*}
			& & \dfrac{283}{3}=90+\left(\dfrac{x^2+5-1}{2\cdot \left(x^2+5\right)-1-4x}  \right)\cdot 5\\
			&\Leftrightarrow& 11x^2-52x+57=0\\ 
			&\Leftrightarrow&  \hoac{&{x=3} \\
			&{x=\dfrac{19}{11}}.} 
		\end{eqnarray*}
		Do $x$  nguyên dương nên suy ra $x=3$.
	 }
\end{ex}

% nguồn: Ok_TracNghiemLop11_Lan1
\begin{ex}%[1D5B1-4]%[ma:658f67e8]
	Người ta ghi lại tuổi thọ của một số con ong cho kết  quả như sau:
	\begin{center}
		\begin{tabular}{|l|c|c|c|c|c|c|}
			\hline Tuổi thọ (ngày) &{$[0;20)$}&{$[20;40)$}&{$[40;60)$}&{$[60;80)$}&{$[80;100)$}\\ 
			\hline Số lượng & $5$ & $12$ & $23$  & $31$ & $29$ \\
			\hline 
		\end{tabular}
	\end{center}
	Nhóm chứa mốt của mẫu số liệu  này là
	\choice 
		{  $[20;40)$}
		 {$[40;60)$}
		 {\True $[60;80)$}
		 {$[80;100)$}
	\loigiai{ 
		Nhóm chứa mốt của mẫu số  liệu này là $[60;80)$.
	} 
\end{ex}


%%%%%%%%%%%%%%%% ĐÚNG/SAI %%%%%%%%%%%%%%%%

% nguồn: Ok_dungsai_lop11_Lan1
\begin{ex}%[1D5N1-4]%[ma:5774ebf4]
 Số lượng người đi xem một bộ phim mới  theo độ tuổi trong một rạp chiếu phim (sau $1$ giờ đầu công chiếu) được ghi lại theo bảng phân phối ghép nhóm sau
 \begin{center}
  \begin{tabular}{|c|c|c|c|c|c|}
   \hline 
   Độ tuổi & $[10;20)$ & $[20;30)$ & $[30;40)$ & $[40;50)$ &  $[50;60)$ \\
    \hline
   Số người & $6$ & $12$ & $16$ & $7$  & $2$ \\
   \hline 
  \end{tabular}
 \end{center}
 \choiceTF 
  {\True Giá trị đại diện nhóm $[50;60)$ là $\dfrac{50+60}{2}=55$} 
  {\True Độ tuổi được dự báo  là ít xem phim đó nhất là thuộc nhóm $[50;60)$}
  {\True  Nhóm chứa mốt là nửa khoảng $[30;40)$}
   {Độ tuổi được dự báo là thích xem phim đó nhiều nhất là $31$ tuổi}
  \loigiai{
  \begin{itemchoice}
    \itemch Đúng. Giá trị đại diện nhóm $[50;60)$ là $55$
   \itemch Đúng. Độ tuổi được dự báo là ít xem phim đó nhất là  thuộc nhóm $[50;60)$ với $2$ người.
   \itemch Đúng. Nhóm chứa mốt là nửa khoảng $[30;40)$. 
   \itemch Sai. Khi đó: $u_m = 30$, $n_m=16$, $n_{m-1}=12$,  $n_{m+1}=7$, $u_{m+1}-u_m =40-30=10$.\\
   Ta có mốt là $M_0 = 30 + \dfrac{16-12}{(16-12)+(16-7)} \cdot 10  = \dfrac{430}{13} \approx 33{,}08$.\\
   Vậy độ tuổi được dự báo là thích xem phim nhiều nhất là $33$  tuổi.
  \end{itemchoice}
  }
\end{ex}

% nguồn: Ok_dungsai_lop11_Lan1
\begin{ex}%[1D5H1-4]%[ma:5ac10c36]
 Chiều cao của 35 cây bạch đàn sinh  trưởng của 12 tháng được ghi lại ở bảng sau (đơn vị: mét).\\
 \centerline{\begin{tabular}{|c|c|c|c|c|c|c|c|c|}
  \hline
 $6{,}6$ & $7{,}5$ & $8{,}2$ & $8{,}2$ & $7{,}8$ & $7{,}9$ & $9{,}0$ &  $8{,}9$ & $8{,}2$\\
 \hline 
 $7{,}2$ & $7{,}5$  & $8{,}3$ & $7{,}4$ & $8{,}7$ & $7{,}7$ & $7{,}0$ & $9{,}4$ & $8{,}7$\\
 \hline 
 $8{,}0$ & $7{,}7$ & $7{,}8$ & $8{,}3$ & $8{,}6$ & $8{,}1$ & $8{,}1$ & $9{,}5$  & $6{,}9$\\
 \hline 
  $8{,}0$ & $7{,}6$ & $7{,}9$ & $7{,}3$ & $8{,}5$ & $8{,}4$ & $8{,}0$ & $8{,}8$ &\\
 \hline 
\end{tabular}}\\
Khi  đó
\choiceTF 
 {\True Chiều cao sinh trưởng trung bình của mỗi tháng của cây bạch đàn  là $8{,}05$ (m)}
 {\True Tổng hợp lại kết quả  chiều cao của cây bạch đàn vào bảng tần số ghép nhóm theo mẫu sau.\\
 \centerline{\begin{tabular}{|c|c|c|c|c|c|c|}
  \hline
 Chiều cao (m) & $[6{,}5;7{,}0)$ & $[7{,}0;7{,}5)$ & $[7{,}5;8)$ & $[8;8{,}5)$ & $[8{,}5;9{,}0)$ &  $[9{,}0;9{,}5)$\\
 \hline 
 Số cây & $2$ & $4$ & $9$ & $11$ & $6$  & $3$\\
  \hline
\end{tabular}}}
{\True Ước lượng chiều cao trung bình mỗi tháng của cây bạch đàn từ bảng tần số ghép nhóm  ở ý (b) là $\approx 8{,}09$ (m)}
{Chiều cao tăng trưởng của cây bạch đàn được $8{,}18$ (m)  là cao nhất}
\loigiai{\begin{itemchoice}
 \itemch Đúng. Chiều cao sinh trưởng trung  bình của mỗi tháng của cây bạch đàn là: $8,05~m$\\
 Chiều cao sinh trưởng  trung bình của mỗi tháng của cây bạch đàn là $8{,}05$ (m).
 \itemch Đúng. Tổng hợp lại kết quả chiều cao của cây bạch đàn vào bảng tần số ghép nhóm theo mẫu sau.\\ 
 \centerline{\begin{tabular}{|c|c|c|c|c|c|c|}
  \hline
 Chiều cao (m) &  $[6{,}5;7{,}0)$ & $[7{,}0;7{,}5)$ & $[7{,}5;8)$ & $[8;8{,}5)$ & $[8{,}5;9{,}0)$ & $[9{,}0;9{,}5)$\\
  \hline
 Số cây & $2$ & $4$  & $9$ & $11$ & $6$ & $3$\\
 \hline 
 \end{tabular}}
 \itemch Đúng. Ước lượng chiều cao trung bình mỗi  tháng của cây bạch đàn từ bảng tần số ghép nhóm ở ý (b) là $\approx 8{,}09$ (m).\\
 \centerline{\begin{tabular}{|c|c|c|c|c|c|c|}
  \hline
 Chiều cao (m) & $[6{,}5;7{,}0)$  & $[7{,}0;7{,}5)$ & $[7{,}5;8)$ & $[8;8{,}5)$ & $[8{,}5;9)$ & $[9;9{,}5)$\\
 \hline 
 Giá trị đại diện & $6{,}75$ &  $7{,}25$ & $7{,}75$ & $8{,}25$ & $8{,}75$ & $9{,}25$\\
  \hline
 Số lần & $2$ & $4$ & $9$ & $11$ & $6$  & $3$\\
 \hline 
 \end{tabular}}\\
 Chiều cao  trung bình mỗi tháng sinh trưởng xấp xỉ bằng
 $$\dfrac{6{,}75\cdot2+7{,}25\cdot4+7{,}75\cdot9+8{,}25\cdot11+8{,}75\cdot6+9{,}25\cdot3}{35}\approx 8{,}09.$$ 
 \itemch Sai. Chiều cao tăng trưởng của cây bạch đàn được $8{,}18$ m là  cao nhất.\\
 Nhóm chứa mốt của mẫu số liệu trên  là $[8;8{,}5)$\\
  Do đó $u_m=8$; ${n_{m-1}}=9$; ${n_{m+1}}=6$; ${u_{m+1}}-u_m=8{,}5-8=0{,}5$.\\
  Mốt của mẫu số liệu ghép nhóm là $\mathrm{M_o}=8+\dfrac{11-9}{(11-9)+(11-6)}\cdot 0{,}5\approx 8{,}14$\\
 Vậy chiều cao tăng trưởng của cây bạch đàn  được $8{,}14$ m là cao nhất.
\end{itemchoice}}
\end{ex}

% nguồn: Ok_dungsai_lop11_Lan1
\begin{ex}%[1D5H1-4]%[ma:c998faaf]
 Người ta tiến hành phỏng vấn $30$ người về một bộ phim mới chiếu trên truyền hình. Người điều tra yêu  cầu cho điểm bộ phim (thang điểm là 100). Kết quả được trình bày trong bảng phân bố tần số ghép nhóm sau đây.\\
 \centerline{\begin{tabular}{|c|c|c|c|c|c|}
 \hline 
 Số điểm &  $[50;60)$ & $[60;70)$ & $[70;80)$ & $[80;90)$ & $[90;100)$\\
  \hline
 Số người & $2$ &  $6$ & $10$ & $8$ & $4$\\
 \hline 
\end{tabular}}\\
 Khi đó
\choiceTF 
 {\True Ước lượng số trung bình của mẫu  ghép là $77$}
 {\True Giá trị đại diện của  nhóm $[90;100)$ là $95$}
 {Nhóm chứa mốt của  mẫu số liệu trên là nhóm $[80;90)$}
 {Mốt của mẫu số liệu là  $74{,}67$}
 \loigiai{
 \centerline{\begin{tabular}{|c|c|c|c|c|c|}
  \hline
 Số điểm & $[50;60)$ & $[60;70)$ & $[70;80)$ & $[80;90)$ & $[90;100)$\\ 
  \hline
 Giá trị đại diện & $55$ & $65$ & $75$ & $85$  & $95$\\
 \hline 
 Số người & $2$ & $6$ & $10$ & $8$ & $4$\\ 
 \hline 
 \end{tabular}}\begin{itemchoice}
 \itemch Đúng. Ước lượng số trung bình của mẫu ghép là  $77$.\\
 Ước lượng số trung bình của mẫu ghép là 
 $$\overline{x}=\dfrac{55.2+65.6+75.10+85.8+95.4}{30}=77.$$ 
 \itemch Đúng. Giá  trị đại diện của nhóm $[90;100)$ là 95.
 \itemch Sai. Nhóm chứa  mốt của mẫu số liệu trên là nhóm $[80;90)$.\\
 Nhóm  chứa mốt của mẫu số liệu trên là nhóm $[70;80)$.
  \itemch Sai. Mốt của mẫu số liệu là $74{,}67$.\\
 Do đó, $u_m=70$; $n_{m-1}=6$; $n_m=10$; $n_{m+1}=8$;  $u_{m+1}-u_m=80-70=10$.\\
  $$\mathrm{M_{o}}=70+\dfrac{10-6}{(10-6)+(10-8)}\cdot 10=76{,}67.$$
 \end{itemchoice}
} 
\end{ex}

% nguồn: Ok_dungsai_lop11_Lan1
\begin{ex}%[1D5H1-4]%[ma:26fb1700]
 Một nhà thực vật học đo chiều dài trung bình $74$ lá cây (đơn vị: milimet) và thu được bảng tần số  ghép nhóm như sau:
 \begin{center}
  \begin{tabular}{|c|c|c|}
    \hline
   Nhóm & Giá trị đại diện & Tần số \\ 
    \hline
   $[5{,}45;5{,}85)$ &  $5{,}65$ & $5$ \\
   \hline 
   $[5{,}85;6{,}25)$ & $6{,}05$  & $9$ \\
   \hline 
   $[6{,}25;6{,}65)$ &  $6{,}45$ & $15$ \\
    \hline
   $[6{,}65;7{,}05)$  & $6{,}85$ & $19$ \\
    \hline
    $[7{,}05;7{,}45)$ & $7{,}25$ & $16$ \\
   \hline 
   $[7{,}45;7{,}85)$ &  $7{,}65$ & $8$ \\
    \hline
   $[7{,}85;8{,}25)$ & $8{,}05$  & $2$ \\
    \hline
  \end{tabular}
 \end{center}
  \choiceTF
  {Chiều dài trung bình của $74$  là cây bằng $\approx 6{,}4 \,(mm)$}
  {\True Độ dài  của nhóm là $0{,}4$}
  {Nhóm chứa mốt là  $[7{,}05;7{,}45)$}
  {Mốt của mẫu số liệu  ghép nhóm là $\approx 6{,}65$}
 \loigiai{ 
  \begin{itemchoice}
   \itemch Sai.  Chiều dài trung bình của $74$ lá cây là
    \[
   \overline{x} =  \dfrac{5{,}65 \cdot 5 + 6{,}05 \cdot 9+ 6{,}45 \cdot 15 + 6{,}85 \cdot 19 + 7{,}25 \cdot 16 + 7{,}65 \cdot 8 + 8{,}05 \cdot 2}{74} = \dfrac{5029}{740} \approx 6{,}8 \text{ (mm)}.
   \] 
   \itemch Đúng. Độ dài  nhóm là $0{,}4$.
   \itemch Sai. Nhóm chứa mốt  là $[6{,}65;7{,}05)$.
   \itemch Sai. Ta có $u_m = 6{,}65$; $u_{m+1}=7{,}05$; $u_{m+1}-u_m = 0{,}4$; $n_m = 19$;  $n_{m-1}=15$; $n_{m+1}=16$.\\
   Do đó mốt của mẫu số liệu ghép nhóm  là
   \[ 
   M_0 = 6{,}65 + \dfrac{19-15}{(19-15)+(19-16)} \cdot 0{,}4 \approx 6{,}88. 
   \] 
  \end{itemchoice}
  }
\end{ex}

% nguồn: Ok_dungsai_lop11_Lan1
\begin{ex}%[1D5H1-4]%[ma:7ca49fc1]
 Cho mẫu số liệu điểm môn  Toán của một nhóm học sinh như sau:
 \begin{center}
  \begin{tabular}{|l|c|c|c|c|}
   \hline Điểm & $[6; 7)$ & $[7; 8)$ & $[8;  9)$ & $[9; 10]$ \\
   \hline Số học sinh  & $8$ & $7$ & $10$ & $5$ \\
   \hline 
  \end{tabular}
 \end{center}
 Xét tính  đúng sai của các mệnh đề sau:
 \choiceTF 
  {\True Mẫu số liệu đã cho là  mẫu số liệu ghép nhóm}
  {\True Cỡ mẫu của mẫu số liệu là $30$} 
  {\True Điểm trung bình  của các học sinh là $7{,}9$}
  {Mốt của mẫu số liệu là  $10$}
  \loigiai{
  Ta có bảng sau 
  \begin{center}
   \begin{tabular}{|l|c|c|c|c|}
    \hline Điểm & {$[6; 7)$} &  {$[7; 8)$} & {$[8; 9)$} & {$[9; 10]$} \\
    \hline Giá trị đại diện & $6{,}5$ & $7{,}5$ & $8{,}5$ & $9{,}5$  \\
    \hline Số học sinh  & $8$ & $7$ & $10$ & $5$ \\
    \hline 
   \end{tabular}
  \end{center}
  \begin{itemchoice}
    \itemch Đúng. Vì theo định nghĩa mẫu số liệu ghép nhóm.
   \itemch  Đúng. Vì cỡ mẫu là $8+7+10+5=30$.
   \itemch Đúng. Vì điểm trung bình của các học sinh là $$\overline{x}=\dfrac{6{,}58+7{,}5 \cdot 7+8{,}5 \cdot 10+9{,}5 \cdot 5}{30}=7{,}9.$$ 
   \itemch Sai. Vì nhóm  chứa Mốt là $[8; 9)$.\\
   Mốt của mẫu  số liệu là $M_e=8+\dfrac{10-7}{10-7+10-5}(9-8)=8{,}375$.
  \end{itemchoice}
 } 
\end{ex}

% nguồn: Ok_dungsai_lop11_Lan1
\begin{ex}%[1D5H1-4]%[ma:15a17382]
 Một nhà nghiên cứu ghi lại thời gian (giờ) sử dụng Facebook của 30 học sinh trong 02 tuần. Kết quả thu  được mẫu số liệu như sau.
 $$\begin{matrix}
 21 & 17 & 22 & 18 & 20 & 17 & 15 & 13 &  15 & 20\\
 15  & 12 & 18 & 17 & 25 & 17 & 21 & 15 & 12 & 18\\
 16 & 23 & 14 & 18 &  19 & 13 & 16 & 19 & 18 & 17\\
\end{matrix}$$
 Khi đó
\choiceTF 
 {Số giờ trung bình của học sinh  trong $02$ tuần $16{,}37$ giờ}
 {\True Tổng hợp kết quả thời gian sử dụng Facebook  của học sinh vào bảng tần số ghép nhóm theo mẫu sau.\\
 \centerline{\begin{tabular}{|c|c|c|c|c|c|}
  \hline
  Số giờ & $[12;15)$ & $[15;18)$ & $[18;21)$ & $[21;24)$ & $[24;27)$\\
 \hline 
 Giá trị đại diện & $13{,}5$ & $16{,}5$ &  $18{,}5$ & $21{,}5$ & $24{,}5$\\
 \hline 
 Số học sinh & $5$ & $12$ &  $8$ & $4$ & $1$\\
  \hline
\end{tabular}}}
 {\True Nhóm chứa mốt của mẫu số liệu ý (b) là nhóm $[15;18)$}
{\True Mốt của  mẫu số liệu ý (b) bằng $16{,}91$}
\loigiai{\begin{itemchoice}
 \itemch Sai. Số giờ trung bình của học sinh trong 02 tuần  $16{,}37$ giờ.\\
 Tổng số thời gian sử dụng Facebook của 30 học sinh  là $521$ giờ.\\
 Số giờ trung bình của học sinh  trong 02 tuần $\overline{x}=\dfrac{521}{30}=17{,}37$ giờ.
 \itemch Đúng. Tổng hợp kết quả  thời gian sử dụng Facebook của học sinh vào bảng tần số ghép nhóm theo mẫu sau.\\
 \centerline{\begin{tabular}{|c|c|c|c|c|c|}
 \hline 
 Số giờ & $[12;15)$ & $[15;18)$ &  $[18;21)$ & $[21;24)$ & $[24;27)$\\
  \hline
 Giá trị đại diện &  $13{,}5$ & $16{,}5$ & $18{,}5$ & $21{,}5$ & $24{,}5$\\
 \hline 
 Số học sinh & $5$  & $12$ & $8$ & $4$ & $1$\\
 \hline 
 \end{tabular}}
 \itemch Đúng. Nhóm chứa mốt của mẫu số liệu ý (b) là  nhóm $[15;18)$ .\\
 Nhóm chứa mốt  của mẫu số liệu trên là nhóm $[15;18)$.
 \itemch Đúng. Mốt của  mẫu số liệu ý (b) bằng $16{,}91$.\\
 Do đó,  $u_m=15$; $n_{m-1}=5$; $n_m=12$; $n_{m+1}=8$; $u_{m+1}-u_m=18-15=3$.\\
  $$\mathrm{M_{o}}=15+\dfrac{12-5}{(12-5)+(12-8)}\cdot 3=16{,}91.$$
 \end{itemchoice}
 }
\end{ex}

% nguồn: Ok_dungsai_lop11_Lan1
\begin{ex}%[1D5H1-4]%[ma:170bcf35]
 Thống kê  điểm trung bình môn Toán của $82$ học sinh lớp $11$ cho trong bảng sau
 \begin{center}
  \begin{tabular}{|c|c|c|c|c|c|c|c|}
   \hline 
   Khoảng  điểm &$[6{,}5;7)$ &$[7;7{,}5)$ &$[7{,}5;8)$ &$[8;8{,}5)$ &$[8{,}5;9)$
    &$[9;5{,}5)$ &$[9{,}5;10)$ \\
    \hline
   Tần số &$8$ &$10$ &$16$  &$24$ & $13$ & $7$& $4$ \\
    \hline
  \end{tabular}
 \end{center}
  Các mệnh đề sau đúng hay sai?
 \choiceTF 
  {\True Số  học sinh đạt từ $8$ điểm trở lên là $48$}
  {\True Giá trị đại diện của nhóm $5$  là $8{,}75$}
  {Điểm trung bình  môn Toán của $82$ học sinh lớp $11$ là $8{,}1$}
  {\True Mốt của mẫu  số liệu ghép nhóm trên là $8{,}21$}
 \loigiai{ 
  Bảng tần số ghép nhóm theo giá trị đại diện là 
  \begin{center}
   \begin{tabular}{|c|c|c|c|c|c|c|c|}
     \hline
    Khoảng điểm  &$[6{,}5;7)$ &$[7;7{,}5)$ &$[7{,}5;8)$ &$[8;8{,}5)$ &$[8{,}5;9)$
    &$[9;5{,}5)$ &$[9{,}5;10)$  \\
     \hline
    Giá trị đại  diện &$6{,}75$& $7{,}25$& $7{,}75$& $8{,}25$& $8{,}75$& $9{,}25$& $9{,}75$\\
     \hline
    Tần  số &$8$ &$10$ &$16$ &$24$ & $13$ & $7$& $4$ \\
     \hline
   \end{tabular}
  \end{center}
  \begin{itemchoice}
   \itemch Số học sinh đạt từ $8$ điểm  trở lên là $24+13+7+4=48$.
    \itemch Giá trị đại điện của nhóm $5$ là $\dfrac{8{,}5+9}{2}=8{,}75$.
   \itemch Điểm trung bình môn Toán của một số học sinh  lớp $11$ là
   $$\overline{x} = \dfrac{8\cdot 6{,}75 + 10\cdot 7{,}25 + 16\cdot 7{,}75  + 24\cdot 8{,}25 + 13 \cdot 8{,}75 + 7\cdot 9{,}25 + 4\cdot 9{,}75}{82} \approx 8{,}12.$$
   \itemch Mốt $M_o$  chứa trong nhóm $[8;8{,}5)$.\\
   Do đó $u_m=8;  u_{m+1} = 8{,}5 \Rightarrow u_{m+1} - u_m = 8{,}5-8 = 0{,}5$.\\
   Và $n_{m-1}=16$;  $n_m = 24$; $n_{m+1} = 13$. Khi đó
   $$M_o = 8 + \dfrac{24-16}{(24-16) + (24-13)} \cdot  0{,}5 = 8{,}21.$$
  \end{itemchoice}
  }
\end{ex}

% nguồn: Ok_dungsai_lop11_Lan2
\begin{ex}%[1D5H1-4]%[ma:3b24d494]
	Số lượng người đi xem một bộ phim mới theo độ tuổi trong một rạp chiếu phim (sau $1$ giờ đầu công chiếu) được ghi lại theo bảng phân  phối ghép nhóm sau
	\begin{center}
		\begin{tabular}{|l|c|c|c|c|c|c|c|}
			\hline Độ tuổi &  {$[10;20)$} & {$ [20;30)$} & {$[30;40) $} & {$[40;50)$} & {$[50;60)$}\\
			\hline Số người& $6$ & $12$ & $16$ & $7$&  $2$\\
			\hline 
		\end{tabular}
	\end{center}
	 \choiceTF
		{\True Giá  trị đại diện nhóm $[50;60)$ là $55$}
		{\True Độ tuổi được dự báo là ít xem  phim đó nhất là thuộc nhóm $[50;60)$}
		{\True Nhóm chứa mốt là  nửa khoảng $[30;40)$}
		{Độ tuổi  được dự báo là thích xem phim đó nhiều nhất là $31$ tuổi}
	\loigiai{ 
		\begin{itemchoice}
			\itemch  Dung. Nhóm chứa mốt là nửa khoảng $[30;40)$.
			\itemch Đúng. Độ  tuổi được dự báo là ít xem phim đó nhất là thuộc nhóm $[50;60)$.
			\itemch Đúng. $u_m=30,n_m=16,n_{m-1}=12,n_{m+1}=7,u_{m+1}-u_m=40-30=10$.\\ 
			 Ta có mốt là $M_o=30+\dfrac{16-12}{(16-12)+(16-7)}\cdot 10=\dfrac{430}{13}\approx 33{,}08$.
			\itemch Sai. Độ tuổi được dự báo là thích xem phim  đó nhiều nhất là $33$ tuổi.
		\end{itemchoice}
	 }
\end{ex}

% nguồn: Ok_dungsai_lop11_Lan2
\begin{ex}%[1D5H1-4]%[ma:53ec1d6d]
	Cho mẫu số liệu về thời gian (phút) đi từ nhà đến nơi làm việc của các nhân viên của một công ty như  sau
	\begin{center}
		\begin{tabular}{|c|c|c|c|c|c|c|c|}
			 \hline
			Thời gian & $[15;20)$ & $[20;25)$ & $[25;30)$ & $[30;35)$ & $[35;40)$ & $[40;45)$ & $[45;50)$ \\ 
			 \hline
			Số nhân viên & $6$ & $14$ & $25$ & $37$ & $21$ & $13$  & $9$ \\
			\hline 
		\end{tabular}
	\end{center}
	\choiceTF 
		{\True Mẫu số  liệu được ghép thành $7$ nhóm}
		{Độ  dài mỗi nhóm của mẫu số liệu bằng $15$}
		{\True Số nhân viên đi từ nhà đến nơi làm việc  từ $25$ phút đến dưới $30$ phút là $25$}
		{Nhóm có số nhân  viên nhiều nhất là nhóm $[45;50)$}
	 \loigiai{
		\begin{itemchoice}
			\itemch  \textbf{Đúng.}\\
			Mẫu số liệu được ghép thành $7$ nhóm. 
			 \itemch \textbf{Sai.}\\
			Độ  dài mỗi nhóm của mẫu số liệu bằng $5$.
			 \itemch \textbf{Đúng.}\\
			 Số nhân viên đi từ nhà đến nơi làm việc từ $25$ phút đến dưới $30$ phút là $25$.
			\itemch  \textbf{Sai.}\\
			Nhóm có số nhân viên nhiều nhất  là nhóm $[30;35)$.
		\end{itemchoice}
	} 
\end{ex}

% nguồn: Ok_dungsai_lop11_Lan2
\begin{ex}%[1D5H1-4]%[ma:3a0dba97]
	Một nhà thực vật học  đo chiều dài trung bình của $74$ lá cây (đơn vị: milimét) và thu được bảng tần số ghép nhóm như sau
	\begin{center}
		\begin{tabular}{|l|c|c|c|}
			\hline Nhóm & Giá trị đại diện  & Tần số\\
			 \hline $[5{,}45;5{,}85)$& $5{,}6$ & $5$\\
			\hline $[5{,}85;6{,}25)$& $6{,}05$ &  $9$\\
			\hline $[6{,}25;6{,}65)$& $6{,}45$ &  $15$\\
			 \hline $[6{,}65;7{,}05)$& $6{,}85$ & $19$\\
			\hline $[7{,}05;7{,}45)$& $7{,}25$ &  $16$\\
			\hline $[7{,}45;7{,}85)$& $7{,}65$  & $8$\\
			\hline  $[7{,}85;8{,}25)$& $8{,}05$ & $2$\\
			 \hline
		\end{tabular}
	\end{center}
	\choiceTF 
		{Chiều dài trung bình của $74$  lá bằng $6{,}4$ (mm) (làm tròn kết quả đến hàng phần chục)}
		{\True Độ dài nhóm  là $0{,}4$}
		{Nhóm chứa mốt là  $[7{,}05;7{,}45)$}
		{mốt của mẫu số liệu ghép nhóm bằng $6{,}65$ (làm tròn kết  quả đến hàng phần trăm)}
	 \loigiai{
		\begin{itemchoice}
			 \itemch Sai. Chiều dài trung bình của $74$ lá cây là
			 $$\overline{x}=\dfrac{5{,}65\cdot 5+6{,}05\cdot 9+6{,}45\cdot 15+6{,}85\cdot 19+7{,}25\cdot 16+7{,}65\cdot 8+8{,}05\cdot 2}{74}=\dfrac{5029}{740}\approx 6{,}8\text{(mm)}.$$
			\itemch Đúng. Độ dài nhóm là  $0{,}4$.
			\itemch Sai. Nhóm chứa  mốt là $[6,65;7,05)$.\\
			Trong đó: $u_m=6{,}65;u_{m+1}=7{,}05;u_{m+1}-u_m=0{,}4;n_m=19;n_{m-1}=15;n_{m+1}=16$. 
			\itemch Sai. Do đó  mốt của mẫu số liệu ghép nhóm là $$M_o=6{,}65+\dfrac{19-15}{(19-15)+(19-16)}\cdot 0{,}4\approx 6{,}88.$$
		\end{itemchoice}
	} 
\end{ex}

% nguồn: Ok_dungsai_lop11_Lan2
\begin{ex}%[1D5H1-4]%[ma:b692e91a]
	Một nhà nghiên cứu ghi lại thời  gian (giờ) sử dụng Facbook của $30$ học sinh trong $02$ tuần. Kết quả thu được mẫu số liệu như sau
	\begin{center}
		\begin{tabular}{|l|c|c|c|c|c|c|c|c|c|c|}
			\hline $21$ & $17$&$22$&$18$& $20$& $17$& $15$& $13$& $15$&  $20$\\
			\hline $15$&  $12$&$18$&$17$&$25$&$17$&$21$&$15$&$12$&$18$\\
			 \hline $16$&$23$&$14$&$18$&$19$&$13$&$16$&$19$&$18$&$17$\\
			\hline 
		\end{tabular}
	\end{center}
	\choiceTF 
		{Số  giờ trung bình của học sinh trong $02$ tuần là $16,37$ giờ (kết quả làm tròn đến hàng phần trăm)}
		{\True Tổng hợp kết quả thời gian sử dụng Facbook của học sinh vào bảng tần số  ghép nhóm theo mẫu sau
		\begin{center}
			\begin{tabular}{|l|c|c|c|c|c|c|c|}
				\hline Số giờ & {$[12{,}5;15)$} & {$ [15;18)$} & {$[18;21) $} & {$[21;24)$} & {$[24;27)$}\\ 
				 \hline Giá trị đại diện & $13{,}5$ & $16{,}5$ & $18{,}5$ & $21{,}5$& $24{,}5$\\
				\hline Số học sinh& $5$ & $12$ & $8$ & $4$& $1$\\ 
				\hline 
		\end{tabular}
	\end{center}}
	{\True Khi mẫu số liệu được ghép nhóm thì nhóm chứa mốt là nhóm $[15;18)$} 
	{\True Khi mẫu số liệu được ghép nhóm thì mốt bằng  $16{,}91$ (kết quả làm tròn đến hàng phần trăm)}
	\loigiai{ 
		\begin{itemchoice}
			\itemch Sai. Tổng số thời gian sử dụng Facbook  của $30$ học sinh là $521$ giờ.\\
			Số giờ trung bình của  học sinh trong $02$ tuần là $\overline{x}=\dfrac{521}{30}\approx17{,}37$ giờ.
			\itemch Đúng. Tổng hợp kết quả thời gian sử dụng Facbook của học sinh  vào bảng tần số ghép nhóm theo mẫu sau
			\begin{center}
				\begin{tabular}{|l|c|c|c|c|c|c|c|}
					\hline Số giờ &  {$[12{,}5;15)$} & {$ [15;18)$} & {$[18;21) $} & {$[21;24)$} & {$[24;27)$}\\
					\hline Số học sinh &  $13{,}5$ & $16{,}5$ & $18{,}5$ & $21{,}5$& $24{,}5$\\
					\hline Giá trị đại diện& $5$ & $12$ &  $8$ & $4$& $1$\\
					 \hline
				\end{tabular}
			\end{center}
			\itemch Đúng. Nhóm chứa mốt của mẫu số liệu trên là nhóm $[15;18)$. 
			\itemch  Đúng. $u_m=15;n_{m-1}=5;n_m=12;n_{m+1}=8;u_{m+1}-u_m=18-15=3{,}0$.\\
			mốt của mẫu số liệu ghép nhóm  là $M_o=15+\dfrac{12-5}{(12-5)+(12-8)}\cdot 3\approx 16{,}91.$
		\end{itemchoice}
	} 
\end{ex}

% nguồn: Ok_dungsai_lop11_Lan2
\begin{ex}%[1D5H1-4]%[ma:6e79775f]
	Số cuộc điện thoại một người thực hiện mỗi ngày trong $30$  ngày được lựa chọn ngẫu nhiên được thống kê trong bảng sau
	\begin{center}
		\begin{tabular}{|l|c|c|c|c|c|c|c|}
			\hline Số cuộc  gọi & {$[2{,}25;5{,}5)$} & {$ [5{,}5;8{,}5)$} & {$[8{,}5;11{,}5) $} & {$[11{,}5;14{,}5)$} & {$[14{,}5;11{,}5)$}\\
			\hline Số  ngày& $5$ & $13$ & $7$ & $3$& $2$\\
			\hline 
		\end{tabular}
	\end{center}
	 \choiceTF
		{Số cuộc gọi trung bình mỗi  ngày là $8,1$}
		{\True  Nhóm chứa mốt là: $[5{,}5;8{,}5)$}
		{\True mốt của mẫu số  liệu ghép nhóm là $\approx 7{,}21$ (kết quả làm tròn đến hàng phần trăm)}
		{Người đó thực hiện tối đa khoảng $8$ cuộc gọi  mỗi ngày}
	\loigiai{ 
		\begin{itemchoice}
			\itemch Sai. Ta viết lại bảng tần số ghép nhóm  theo giá trị đại diện là
			\begin{center}
				\begin{tabular}{|l|c|c|c|c|c|c|c|}
					\hline Số cuộc gọi & {$[2{,}25;5{,}5)$} & {$ [5{,}5;8{,}5)$} & {$[8{,}5;11{,}5) $} & {$[11{,}5;14{,}5)$} & {$[14{,}5;11{,}5)$}\\ 
					\hline Giá trị đại diện& $4$ & $7$ & $10$ &  $13$& $16$\\
					\hline Số ngày& $5$ & $13$  & $7$ & $3$& $2$\\
					\hline 
				\end{tabular}
			\end{center}
			Số cuộc gọi trung bình mỗi ngày là $$\overline{x}=\dfrac{4\cdot 5+7\cdot 13+10\cdot 7+13\cdot 3+16\cdot 2}{30}=8{,}4.$$ 
			\itemch  Đúng. Nhóm chứa mốt là $[5{,}5;8{,}5)$.\\
			Ta có  $u_m=5{,}5;u_{m+1}=8{,}5\Rightarrow u_{m+1}-u_m=3;n_m=13;n_{m-1}=5;n_{m+1}=7$.
			\itemch  Đúng. mốt của mẫu số liệu ghép nhóm là $M_o=5{,}5+\dfrac{13-5}{(13-5)+(13-7)}\cdot 3=\dfrac{101}{14}\approx 7{,}21.$
			\itemch Sai. Người đó  thực hiện tối đa khoảng $7$ cuộc gọi mỗi ngày.
		\end{itemchoice}
	} 
\end{ex}

% nguồn: Ok_dungsai_lop11_Lan1
\begin{ex}%[1D5V1-4]%[1D5V2-3]%[ma:373d247b]
 Bảng tần số ghép nhóm số liệu dưới đây  thống kê cân nặng của 40 học sinh lớp 11A trong một trường trung học phổ thông (đơn vị: kilôgam).
 \begin{center}
  $\begin{array}{|c|c|c|}
  \hline \text { Nhóm } & \text { Tần số } & \begin{array}{c}
  \text { Tần số  } \\
  \text { tích  luỹ }
 \end{array} \\
 \hline[30; 40)  & 2 & 2 \\
 {[40; 50)} & 10 & 12  \\
 {[50; 60)} & 16 &  28 \\
 {[60; 70)} & 8 & 36 \\ 
 {[70;  80)} & 2 & 38 \\
 {[80; 90)} & 2 & 40 \\ 
 \hline & n=40  & \\
  \hline
\end{array}$
\end{center}
\choiceTF 
 {\True Cân nặng trung bình của $40$ học sinh trên là  $56$ kg}
 {Tứ phân vị thứ hai  của mẫu số liệu trên là $Q_2=50$ kg}
 {Giá trị của hiệu tứ phân vị thứ ba với thứ nhất của mẫu số liệu  trên là $Q_3-Q_1=15$ kg}
 {\True Trong bảng số liệu trên học sinh có cân nặng khoảng $54$ kg chiếm nhiều  nhất}
\loigiai{ 
 \begin{itemchoice}
  \itemch Cân nặng trung bình $$\bar{x}=\dfrac{35\cdot 2+45\cdot 10+55\cdot 16+65\cdot 8+75\cdot 2+85\cdot  2}{40}=56.$$
  \itemch Tứ phân  vị thứ hai của dãy số liệu $x_1; x_2; \ldots; x_{40}$ thuộc $[50; 60)$. Do đó tứ phân vị thứ hai của mẫu số liệu ghép nhóm là
  $$ 
  Q_2=50+\dfrac{\dfrac{2\cdot 40}{4}-(2+10)}{16}  \cdot(60-50)=55.
   $$
  \itemch Tứ phân vị thứ nhất $Q_1=48$  (kg) và tứ phân vị thứ ba $Q_3=62{,}5$ (kg). \\
   Nên $Q_3-Q_1=62{,}5-48=14{,}5$.
  \itemch Theo mẫu số  liệu trên ta có mốt
   $$M_{\circ}=50+\dfrac{16-10}{(16-10)+(16-8)}(65-50)=54{,}286.$$
 \end{itemchoice}
} 
\end{ex}

% nguồn: Ok_dungsai_lop11_Lan2
\begin{ex}%[1D5V1-4]%[ma:e6acfa30]
	Thống kê điểm trung bình môn Toán của $82$ học sinh lớp $11$ được cho ở bảng  sau
	\begin{center}
		\begin{tabular}{|c|c|c|c|c|c|c|c|}
			\hline 
			Khoảng điểm & $[6{,}5; 7)$ & $[7; 7{,}5)$ & $[7{,}5; 8)$ & $[8; 8{,}5)$ & $[8{,}5; 9)$ & $[9; 9{,}5)$ & $[9{,}5;  10)$ \\
			 \hline
			Tần số &  $8$ & $10$ & $16$ & $24$ & $13$ & $7$ & $4$ \\
			\hline 
		\end{tabular}
	\end{center}
	\choiceTF 
		{\True Số học sinh  đạt từ $9$ điểm trở lên là $11$}
		{\True Giá trị đại diện  của nhóm $4$ là $8{,}25$}
		{Điểm trung bình môn Toán của $82$ hoc sinh lớp $11$  là $8{,}3$}
		{\True Mốt của mẫu số liệu là  $8{,}21$}
	\loigiai{ 
		Bảng tần số ghép nhóm theo giá trị đại diện  là
		\begin{center}
			\begin{tabular}{|c|c|c|c|c|c|c|c|}
				\hline 
				Khoảng  điểm & $[6,5; 7)$ & $[7; 7,5)$ & $[7,5; 8)$ & $[8; 8,5)$ & $[8,5; 9)$ & $[9; 9,5)$ & $[9,5; 10)$ \\
				 \hline
				Giá trị đại diện & $6{,}75$ & $7{,}25$ & $7{,}75$ & $8{,}25$ & $8{,}75$ &  $9{,}25$ & $9{,}75$ \\
				\hline 
				Tần số & $8$ & $10$ & $16$ & $24$ & $13$ & $7$  & $4$ \\
				\hline 
			\end{tabular}
		\end{center}
		\begin{itemchoice}
			 \itemch \textbf{Đúng.}\\
			Số học sinh đạt từ $9$ điểm trở lên là $11$. 
			\itemch \textbf{Sai.}\\ 
			Giá trị đại diện của  nhóm $4$ là $8{,}25$.
			 \itemch \textbf{Đúng.}\\
			 Điểm trung bình môn Toán của một số hoc sinh lớp $11$ là
			\[\overline{x}=\dfrac{8\cdot 6{,}75+10\cdot 7{,}25+16\cdot 7{,}75+24\cdot 8{,}25+13\cdot 8{,}75+7\cdot  9{,}25+4\cdot 9{,}75}{82} \approx 8{,}12.\]
			\itemch  \textbf{Đúng.}\\
			Mốt $M_{o}$ chứa trong nhóm $[8;8{,}5)$. Do  đó $u_{m}=8$; $u_{m+1}=8{,}5$ suy ra $u_{m+1}-u_{m}=8{,}5-8=0{,}5$; $n_{m-1}=16$; $n_{m}=24$; $n_{m+1}=13$. Ta có
			\[M_{o}=8+\dfrac{24-16}{(24-16)+(24-13)}\cdot(8{,}5-8)=8{,}21.\] 
		\end{itemchoice}
	} 
\end{ex}

% nguồn: Ok_dungsai_lop11_Lan2
\begin{ex}%[1D5V1-4]%[ma:2eb3b62f]
	Cho bảng tần số về  cân nặng của $180$ người dân trong một xã như sau (đơn vị: kg)
	\begin{center}
		\begin{tabular}{|c|c|}
			 \hline
			 Nhóm & Tần số \\
			 \hline
			$[0;10)$ &  $6$ \\
			 \hline
			$[10;20)$ &  $15$ \\
			 \hline
			 $[20;30)$ & $37$ \\
			 \hline
			$[30;40)$ & $48$  \\
			\hline 
			$[40;50)$ & $22$ \\ 
			 \hline
			$[50;60)$  & $29$ \\
			 \hline
			$[60;70)$ &  $23$ \\
			\hline 
		\end{tabular}
	\end{center}
	 \choiceTF
		{\True  Giá trị đại diện của nhóm $2$ là $15$}
		 {Có $100$ người có cân nặng nhỏ hơn $40$ kg}
		{\True Cân nặng trung  bình của người dân là $38{,}5$ kg}
		{\True Mốt của mẫu số liệu là $32{,}97$  kg}
	\loigiai{ 
		Bảng tần số ghép nhóm theo giá  trị đại diện là
		\begin{center}
			\begin{tabular}{|c|c|c|c|} \hline
				Nhóm & Tần số & Giá trị đại diện  & Tần số tích lũy \\ \hline
				$[0;10)$ & $6$ & $5$  & $6$ \\ \hline
				$[10;20)$ &  $15$ & $15$ & $21$ \\ \hline
				$[20;30)$ &  $37$ & $25$ & $58$ \\ \hline
				$[30;40)$  & $48$ & $35$ & $106$ \\ \hline
				$[40;50)$ & $22$  & $45$ & $128$ \\ \hline
				$[50;60)$ &  $29$ & $55$ & $157$ \\ \hline
				$[60;70)$ & $23$ & $65$ & $180$ \\  \hline
			\end{tabular}
		\end{center}
		\begin{itemchoice}
			 \itemch \textbf{Đúng.}\\
			Giá trị đại diện của  nhóm $2$ là $\dfrac{10+20}{2}=15$.
			\itemch \textbf{Sai.}\\ 
			Số  có cân nặng nhỏ hơn $40$ kg là $106$ người.
			\itemch  \textbf{Đúng.}\\
			Cân nặng trung  bình của người dân là
			\[\dfrac{5\cdot6+15\cdot15+25\cdot37+35\cdot48+45\cdot22+55\cdot29+65\cdot23}{180}=38{,}5.\] 
			\itemch \textbf{Đúng.}\\ 
			Nhóm có tần số lớn nhất là nhóm $4$. Nhóm $4$ có đầu mút trái $u=30$; độ dài $g=10$; tần số $n_{4}=48$. Nhóm $3$ và nhóm $5$ có tần  số tương ứng là $n_{3}=37$; $n_{5}=22$. Ta có
			 \[M_{o}=u+(\dfrac{n_{4}-n_{3}}{2n_{4}-n_{3}-n_{5}})\cdot g=30+\dfrac{48-37}{2\cdot 48-37-22} \cdot 10=32{,}97.\]
		\end{itemchoice}
	 }
\end{ex}

% nguồn: Ok_dungsai_lop11_Lan2
\begin{ex}%[1D5V1-4]%[ma:2bbfd598]
	Thống kê điểm trung bình môn Toán của một số  hoc sinh lớp $11$ được cho ở bảng sau
	\begin{center}
		\begin{tabular}{|c|c|c|c|c|c|c|c|}
			 \hline
			Khoảng điểm & $[6{,}5; 7)$ & $[7; 7{,}5)$ & $[7{,}5; 8)$ & $[8; 8{,}5)$ & $[8{,}5; 9)$ & $[9; 9{,}5)$ & $[9{,}5; 10)$ \\ 
			 \hline
			Tần số & $8$ & $10$ & $16$ & $24$ & $13$ &  $7$ & $4$\\
			\hline 
		\end{tabular}
	\end{center}
	 \choiceTF
		{\True Tần số của nhóm $[9{,}5;10)$ là $4$} 
		{Giá trị đại diện của  nhóm $[7;7{,}5)$ là $7$}
		{\True Điểm trung bình môn Toán của một số học sinh lớp $11$  là $8{,}12$}
		{Mốt của mẫu số liệu ghép nhóm trên $8{,}12$} 
	\loigiai{ 
		Bảng tần số ghép nhóm theo giá trị đại diện là 
		\begin{center}
			\begin{tabular}{|c|c|c|c|c|c|c|c|}
				 \hline
				 Khoảng điểm & $[6{,}5; 7)$ & $[7; 7{,}5)$ & $[7{,}5; 8)$ & $[8; 8{,}5)$ & $[8{,}5; 9)$ & $[9; 9{,}5)$ & $[9{,}5; 10)$ \\
				\hline 
				Giá trị đại diện & $6{,}75$ & $7{,}25$ & $7{,}75$ & $8{,}25$ &  $8{,}75$ & $9{,}25$ & $9{,}75$ \\
				\hline 
				 Tần số & $8$ & $10$ & $16$ & $24$ & $13$ & $7$ & $4$\\
				 \hline
			\end{tabular}
		\end{center}
		\begin{itemchoice}
			\itemch \textbf{Đúng.}\\ 
			Tần  số của nhóm $[9{,}5;10)$ là $4$.
			 \itemch \textbf{Sai.}\\
			Giá trị đại diện của nhóm $[7;7{,}5)$là $\dfrac{7+7{,}5}{2}=7{,}25$. 
			\itemch  \textbf{Đúng.}\\
			Điểm trung bình môn Toán  của một số hoc sinh lớp $11$ là
			\[\overline{x}=\dfrac{8\cdot 6{,}75+10\cdot 7{,}25+16\cdot 7{,}75+24\cdot 8{,}25+13\cdot 8{,}75+7\cdot 9{,}25+4\cdot 9{,}75}{82} \approx 8{,}12.\] 
			\itemch \textbf{Sai.}\\ 
			Mốt $M_{o}$ chứa  trong nhóm $[8;8,5)$. Do đó $u_{m}=8$; $u_{m+1}=8{,}5$ suy ra $u_{m+1}-u_{m}=8{,}5-8=0{,}5$; $n_{m-1}=16$; $n_{m}=24$; $n_{m+1}=13$.
			\[M_{o}=8+\dfrac{24-16}{(24-16)+(24-13)}\cdot(8{,}5-8)=8{,}21.\] 
		\end{itemchoice}
	} 
\end{ex}


%%%%%%%%%%%%%%%% TRẢ LỜI NGẮN %%%%%%%%%%%%%%%%

% nguồn: Ok_TraLoiNgan_Lop11_Lan1
\begin{ex}%[1D5N1-4]%[ma:d558fa1a]
	Người ta tiến hành phỏng vấn $50$ người về phim chiếu rạp Lật mặt  $6$ của Lý Hải. Người điều tra yêu cầu cho điểm phim theo thang điểm $100$. Kết quả được trình bày trong bảng phân bố tần số ghép lớp sau đây:
	\begin{center}
		\begin{tabular}{|c|c|c|c|c|c|}
			 \hline
			Số  điểm&$[50;60)$&$[60;70)$&$[70;80)$&$[80;90)$&$[90;100)$\\
			\hline 
			Số  người&$4$&$7$&$9$&$18$&$12$\\
			 \hline
		\end{tabular}
	\end{center}
	Hãy ước lượng mốt của mẫu số  liệu ghép nhóm ở trên.
	\shortans{80} 
	\loigiai{ 
		\begin{center}
			\begin{tabular}{|c|c|c|c|c|c|}
				 \hline
				Số điểm&$[50;60)$&$[60;70)$&$[70;80)$&$[80;90)$&$[90;100)$\\ 
				 \hline
				 Giá trị đại diện&$55$&$65$&$75$&$85$&$95$\\
				\hline 
				Số người&$4$&$7$&$9$&$18$&$12$\\ 
				\hline 
			\end{tabular}
		\end{center}
		Nhóm  chứa mốt của mẫu số liệu trên là $[80;90)$.\\
		 Do đó: $u_m=80;{n_{m-1}}=9;{n_{m+1}}=12;{u_{m+1}}-u_m=90-80=10$.\\
		Mốt của mẫu số liệu  ghép nhóm là \[M_o=80+\dfrac{18-9}{(18-9)+(18-12)}\cdot 10=80.\]
	 }
\end{ex}

% nguồn: Ok_TraLoiNgan_Lop11_Lan1
\begin{ex}%[1D5N1-4]%[ma:89ea613c]
	Kết quả đo chiều cao của $250$ cây dừa đột biến $3$ năm tuổi ở một viện nghiên cứu được tổng hợp ở bảng  sau:\\
	\begin{center}
		\begin{tabular}{|c|c|c|c|c|c|}
			 \hline
			Chiều  cao (m$^2$)&$[8{,}5;8{,}8)$&$[8{,}8;9{,}1)$&$[9{,}1;9{,}4)$&$[9{,}4;9{,}7)$&$[9{,}7;10)$\\
			 \hline
			Số  cây&$36$&$45$&$83$&$65$&$21$\\
			 \hline
		\end{tabular}
	\end{center}
	 Tính mốt của mẫu số liệu ghép nhóm trên (kết quả làm tròn đến hàng phần mười).
	 \shortans{9,3}
	 \loigiai{
		\begin{center}
			\begin{tabular}{|c|c|c|c|c|c|}
				 \hline
				Chiều cao (m$^2$)&$[8{,}5;8{,}8)$&$[8{,}8;9{,}1)$&$[9{,}1;9{,}4)$&$[9{,}4;9{,}7)$&$[9{,}7;10)$\\ 
				\hline 
				 Giá trị đại diện&$8{,}65$&$8{,}95$&$9{,}25$&$9{,}55$&$9{,}85$\\
				 \hline
				Số  cây&$36$&$45$&$83$&$65$&$21$\\
				\hline 
			\end{tabular}
		\end{center}
		Nhóm chứa mốt của  mẫu số liệu trên là $[9{,}1;9{,}4)$.\\
		Do đó $u_m=9{,}1;{n_{m-1}}=45;{n_{m+1}}=65;{u_{m+1}}-u_m=9{,}4-9{,}1=0{,}3$.\\ 
		Mốt của mẫu số liệu ghép nhóm  là \[M_o=9{,}1+\dfrac{83-45}{(83-45)+(83-65)}\cdot 0{,}3\approx 9{,}3.\]
	} 
\end{ex}

% nguồn: Ok_TraLoiNgan_Lop11_Lan2
\begin{ex}%[1D5N1-4]%[ma:4ab7c4fc]
	Khảo sát thời gian tập thể dục của một số học sinh khối $11$ thu được mẫu số liệu ghép nhóm sau: 
	\begin{center}
		\begin{tabular}{|c|c|c|c|c|c|}
			\hline 
			Thời  gian (phút) & [0; 20) & [20; 40) & [40; 60) & [60; 80) & [80; 100) \\
			\hline 
			 Số học sinh & 5 & 9 & 12 & 10 & 6 \\
			 \hline
		\end{tabular}
	\end{center}
	Nhóm chứa mốt của mẫu số liệu trên là nhóm $[a;b)$.  Tính $a+b$
	\shortans{100} 
	 \loigiai{
		Mốt $  M_{o}$ thuộc nhóm $[40;60)$. Nên $a+b=40+60=100$
	 }
\end{ex}

% nguồn: Ok_TraLoiNgan_Lop11_Lan2
\begin{ex}%[1D5N1-4]%[ma:c3bfbe53]
	Mức thưởng  tết cho các nhân viên của một công ty được thống kê trong bảng sau:
	\begin{center}
		\begin{tabular}{|c|c|c|c|c|c|}
			 \hline
			Mức thưởng tết & [5;  10) & [10; 15) & [15; 20) & [20; 25) & [25; 30) \\ \hline
			Số nhân viên & 13 & 35 &  47 & 25 & 10 \\ \hline
		\end{tabular}
	\end{center}
	Số mốt của mẫu số liệu ghép nhóm  trên là
	\shortans{1} 
	\loigiai{ 
		Tần số lớn nhất của mẫu số liệu ghép nhóm trên là $47$. Nên nhóm chứa mốt của mẫu số liệu trên là nhóm $[15; 20)$. Do  đó mẫu số liệu này chỉ có $1$ mốt.
	 }
\end{ex}

% nguồn: Ok_TraLoiNgan_Lop11_Lan2
\begin{ex}%[1D5N1-4]%[ma:3b2ad325]
	Thời gian truy cập internet mỗi buổi tối của  một số học sinh cho trong bảng sau:
	\begin{center}
		\begin{tabular}{|c|c|c|c|c|c|}
			 \hline
			Thời gian (phút) & $[9{,}5; 12{,}5)$ &$ [12{,}5; 15{,}5)$ &$ [15{,}5; 18{,}5)$ &  $[18{,}5; 21{,}5)$ & $[21{,}5; 24{,}5)$ \\
			\hline 
			Số học sinh & 3 & 12 & 15 & 24 & 2  \\
			\hline 
		\end{tabular}
	\end{center}
	Nhóm chứa mốt của mẫu  số liệu trên là nhóm $[a;b)$. Tính $a+b$
	 \shortans{40}
	\loigiai{ 
		Mốt  $ M_{o}$ chứa trong nhóm $[18{,}5; 21{,}5)$. Nên $a+b=18{,}5+21{,}5=40$
	 }
\end{ex}

% nguồn: Ok_TraLoiNgan_Lop11_Lan2
\begin{ex}%[1D5N1-4]%[ma:8213533d]
	Khảo sát thời gian tập thể dục  của một số học sinh khối $11$ thu được mẫu số liệu ghép nhóm sau:
	\begin{center}
		\begin{tabular}{|c|c|c|c|c|c|}
			 \hline
			Thời gian (phút) & [0; 20) &  [20; 40) & [40; 60) & [60; 80) & [80; 100) \\ \hline
			Số học sinh & 5 & 9 & 12 & 10 & 6  \\ \hline
		\end{tabular}
	\end{center}
	Nhóm chứa mốt của mẫu số  liệu trên là nhóm $[a;b)$. Tính $2a+b$
	 \shortans{140}
	 \loigiai{
		Tần số của nhóm $[40; 60)$ lớn nhất nên mốt thuộc  nhóm $[40; 60)$. Vậy $2a+b=2\cdot 40+60=140$
	} 
\end{ex}

% nguồn: Ok_TraLoiNgan_Lop11_Lan2
\begin{ex}%[1D5N1-4]%[ma:4a7ff066]
	Thời gian truy cập internet mỗi buổi tối của một số học sinh cho  trong bảng sau:
	\begin{center}
		\begin{tabular}{|c|c|c|c|c|c|}
			\hline 
			Thời gian (phút) & [9,5; 12,5) & [12,5; 15,5) &  [15,5; 18,5) & [18,5; 21,5) & [21,5; 24,5) \\
			 \hline
			Số học sinh & 3 & 12  & 15 & 24 & 2 \\
			\hline 
		\end{tabular}
	\end{center}
	Nhóm chứa mốt của mẫu số liệu trên là nhóm $[a;b)$.  Tính $a+3b$
	 \shortans{83}
	\loigiai{ 
		Mốt  $ M_{o}$ thuộc nhóm $[18{,}5; 21{,}5)$. Vậy $a+3b=18{,}5+3\cdot 21{,}5=83$
	 }
\end{ex}

% nguồn: Ok_TraLoiNgan_Lop11_Lan2
\begin{ex}%[1D5N1-4]%[ma:e0836b05]
	Khảo sát thời gian tự học bài ở nhà của một  số em học sinh lớp $11$ thu được mẫu ghép nhóm số lượng như sau:
	\begin{center}
		\begin{tabular}{|c|c|c|c|c|c|}
			\hline 
			Thời gian (phút) & [0; 30) & [30; 60) & [60; 90) &  [90; 120) & [120; 150) \\ \hline
			Số học sinh & 9 & 10 & 9 &  15 & 7 \\ \hline
		\end{tabular}
	\end{center}
	Nhóm chứa mốt  là nhóm $[a;b)$. Tính $a+b$
	\shortans{210} 
	\loigiai{ 
		Nhóm  chứa mốt của mẫu số liệu ghép nhóm là nhóm $[90; 120)$. Vậy $a+b=90+120=210$
	 }
\end{ex}

% nguồn: Ok_TraLoiNgan_Lop11_Lan2
\begin{ex}%[1D5N1-4]%[ma:5fde8547]
	Bảng sau cho ta cân nặng của  học sinh một lớp $11$
	\begin{center}
		\begin{tabular}{|l|c|c|c|c|c|c|c|}
			\hline  Cân nặng (kg) & {$[40{,}5;45{,}5)$} & {$ [45{,}5;50{,}5)$} & {$[50{,}5;55{,}5) $} & {$[55{,}5;60{,}5)$} & {$[60{,}5;65{,}5)$}\\
			\hline Số học sinh  & $10$ & $7$ & $16$ & $4$& $2$\\
			 \hline
		\end{tabular}
	\end{center}
	Tìm cân nặng trung bình của học  sinh lớp $11$ đó (kết quả làm tròn đến hàng phần trăm).
	\shortans{50,6} 
	 \loigiai{
		Bảng thống kê cân nặng của học sinh  theo giá trị đại diện
		\begin{center}
			\begin{tabular}{|l|c|c|c|c|c|c|c|}
				\hline Cân nặng (kg) &  {$[40{,}5;45{,}5)$} & {$ [45{,}5;50{,}5)$} & {$[50{,}5;55{,}5) $} & {$[55{,}5;60{,}5)$} & {$[60{,}5;65{,}5)$}\\
				\hline Giá trị đại diện & $43$ & $48$ & $83$ & $65$& $21$\\ 
				\hline Số học sinh & $10$  & $7$ & $16$ & $4$& $2$\\
				 \hline
			\end{tabular}
		\end{center}
		Cân nặng trung bình của học  sinh lớp $11$ đó là
		$$\overline{x}=\dfrac{43\cdot 10+48\cdot 7+53\cdot 16+58\cdot 4+63\cdot 2}{39}\approx  50{,}56\text{ (kg)}.$$
	 }
\end{ex}

% nguồn: Ok_TraLoiNgan_Lop11_Lan2
\begin{ex}%[1D5N1-4]%[ma:acbe6bf0]
	Khảo sát thời gian tập thể dục trong ngày của một  số học sinh khối $11$ thu được mẫu số liệu ghép nhóm sau:
	\begin{center}
		\begin{tabular}{|c|c|c|c|c|c|}
			 \hline
			Thời gian (phút) & [0; 20) & [20; 40) & [40; 60) & [60; 80) & [80; 100) \\ 
			 \hline
			Số học sinh & 5  & 9 & 12 & 10 & 6 \\
			 \hline
		\end{tabular}
	\end{center}
	Số mốt của mẫu số liệu ghép nhóm  trên là
	\shortans{1} 
	\loigiai{ 
		Vì đây là mẫu số liệu  ghép nhóm và tần số các nhóm khác nhau nên có 1 mốt.
	} 
\end{ex}

% nguồn: Ok_TraLoiNgan_Lop11_Lan2
\begin{ex}%[1D5N1-4]%[ma:ca7c9f2b]
	Khảo sát thời gian tập thể dục của một số học sinh khối $11$ thu được mẫu số liệu ghép nhóm  sau:
	\begin{center}
		\begin{tabular}{|c|c|c|c|c|c|}
			\hline 
			Thời gian (phút) & [0; 20) & [20; 40) & [40; 60) & [60; 80)  & [80; 100) \\ \hline
			Số học sinh & 5 & 9  & 12 & 10 & 6 \\ \hline
		\end{tabular}
	\end{center}
	Nhóm chứa mốt của mẫu số  liệu trên là nhóm $[a;b)$. Tính $3a+2b$
	\shortans{240} 
	\loigiai{ 
		Nhóm  chứa mốt của mẫu số liệu ghép nhóm là nhóm có tần số lớn nhất.
		Do đó nhóm chứa mốt của mẫu số liệu ghép nhóm trên là nhóm $[40; 60)$. Vậy $3a+2b=3\cdot 40+2\cdot 60=240$ 
	 }
\end{ex}

% nguồn: Ok_TraLoiNgan_Lop11_Lan2
\begin{ex}%[1D5N1-4]%[ma:26549066]
	Khi thống kê chiều cao  của học sinh khối $12$ trong một trường trung học, ta thu được mẫu số liệu ghép nhóm sau:
	\begin{center}
		\begin{tabular}{|c|c|c|c|c|c|c|}
			 \hline
			Chiều cao & [150; 156) & [156; 162) & [162; 168) & [168; 174) & [174; 180)  & [180; 186) \\ \hline
			Số học sinh & 5 &  18 & 40 & 26 & 8 & 3 \\ \hline
		\end{tabular}
	\end{center}
	 Nhóm chứa mốt của mẫu số liệu trên là nhóm $[a;b)$. Tính $a-2b$
	 \shortans{-174}
	\loigiai{ 
		Nhóm chứa mốt của  mẫu số liệu ghép nhóm là nhóm có tần số lớn nhất.
		 Do đó nhóm chứa mốt của mẫu số liệu ghép nhóm trên là nhóm $[162; 168)$. Vậy $a-2b=162-2\cdot 168=-174$
	 }
\end{ex}

% nguồn: Ok_TraLoiNgan_Lop11_Lan1
\begin{ex}%[1D5H1-4]%[ma:377977f5]
	Bảng sau thống kê số ca nhiễm mới  SARS-CoV-2 mỗi ngày trong tháng 12/2021 tại Việt Nam.
	\begin{center}
		\begin{tabular}{|c|c|c|c|c|c|c|c|} \hline
			{Ngày} & {Số ca} & {Ngày} & {Số ca} & {Ngày} & {Số ca} & {Ngày} & \text{Số  ca} \\ \hline
			1 & 15\,139 & 9 & 15\,965 & 17 &  15\,685 & 25 & 16\,046 \\\hline
			2 &  14\,295 & 10 & 15\,474 & 18 & 16\,363 & 26 & 15\,667 \\\hline
			3 & 14\,254 & 11 & 16\,830 & 19 & 16\,586 & 27 &  15\,310 \\\hline
			4 & 14\,598 & 12 & 15\,264 &  20 & 15\,420 & 28 & 14\,866 \\\hline
			5 & 14\,927 &  13 & 16\,035 & 21 & 16\,806 & 29 & 14\,299 \\\hline
			6 & 15\,215 & 14 & 15\,871 & 22 & 17\,044 & 30 & 20\,454 \\\hline 
			7 & 14\,433 & 15 & 16\,192 & 23 & 16\,860 &  31 & 17\,004 \\\hline
			8 & 15\,223 & 16 & 15\,720 & 24 & 16\,633  &&\\ \hline
		\end{tabular}
		\begin{flushright}
			(\textit{Nguồn}: worldmeters.info) 
		\end{flushright}
	\end{center}
	Tìm mốt của  mẫu số liệu ghép nhóm trên.
	\shortans{15,7} 
	 \loigiai{
		\textbf{\color{magenta}Bảng tần  số ghép nhóm}
		\begin{center}
			\begin{tabular}{|c|c|c|c|c|c|} \hline
				Số ca  (nghìn) & $\left[14;15{,}5\right)$ & $\left[15{,}5;17\right)$ & $\left[17;18{,}5\right)$ & $\left[18{,}5;20\right)$ & $\left[20;21{,}5\right)$ \\ \hline
				Số  ngày & 13 & 15 & 2 & 0 & 1\\ \hline
			\end{tabular}
		\end{center}
		Mốt $M_o$ chứa  trong nhóm $\left[u_{m};u_{m+1}\right)$
		$$M_o=u_{m}+\dfrac{n_{m}-n_{m-1}}{\left(n_{m}-n_{m-1}\right)+\left(n_{m}-n_{m+1}\right)}\left(u_{m+1}-u_{m}\right)$$ 
		Nhóm chứa  mốt của mẫu số liệu trên là nhóm $\left[120;175\right)$\\
		Do đó $\begin{aligned}
		&  u_{m}=15{,}5; u_{m+1}=17 \Rightarrow u_{m+1}-u_{m}=17-15{,}5=1{,}5.\\
		& n_{m-1}=13;  n_{m}=15; n_{m+1}=2. \\
		& M_o=15{,}5+\dfrac{15-13}{(15-13)+(15-2)}\cdot(17-15{,}5) =\dfrac{157}{10} =  15{,}7.
	\end{aligned}$
} 
\end{ex}

% nguồn: Ok_TraLoiNgan_Lop11_Lan1
\begin{ex}%[1D5H1-4]%[ma:c3ca5afd]
	Số cuộc gọi điện thoại một nguời thực hiện mỗi ngày trong $30$ ngày được lựa chọn  ngẫu nhiên được thống kê trong bảng sau:
	\begin{center}
		\begin{tabular}{|c|c|c|c|c|c|}
			\hline Số cuộc gọi &{$[3; 5]$} &{$[6; 8]$}  &{$[9; 11]$} &{$[12; 14]$} &{$[15; 17]$} \\
			\hline Số ngày & 5 & 13 &  7 & 3 & 2 \\
			 \hline
		\end{tabular}
	\end{center}
	Tìm  mốt của mẫu số liệu ghép nhóm trên (làm tròn đến hàng phần mười).
	 \shortans{7,2}
	\loigiai{ 
		Do số cuộc gọi là số nguyên nên ta hiệu chỉnh lại như sau: 
		\begin{center}
			\begin{tabular}{|c|c|c|c|c|c|}
				\hline Số cuộc gọi &{$[2{,}5; 5{,}5)$} &{$[5{,}5;  8{,}5)$} &{$[8{,}5; 11{,}5)$} &{$[11{,}5; 14{,}5)$} &{$[14{,}5; 17{,}5)$} \\
				\hline Số ngày & 5 & 13 & 7 & 3  & 2 \\
				 \hline
			\end{tabular}
		\end{center}
		Nhóm chứa mốt của mẫu  số liệu trên là nhóm $[5{,}5; 8{,}5)$.\\
		Do đó  $u_m=5{,}5$; $n_{m-1}=5$; $n_m=13$; $n_{m+1}=7$; $u_{m+1}-u_m=8{,}5-5{,}5=3$.\\
		Mốt  của mẫu số liệu ghép nhóm là
		\[M_0=5{,}5+\dfrac{13-5}{(13-5)+(13-7)} \cdot 3=\dfrac{101}{14} \approx  7{,}2. \]
	} 
\end{ex}

% nguồn: Ok_TraLoiNgan_Lop11_Lan1
\begin{ex}%[1D5H1-4]%[ma:a96dd9bb]
	Người ta tiến hành phỏng vấn $50$ người về phim chiếu rạp Lật mặt $6$ của Lý Hải. Người điều  tra yêu cầu cho điểm phim theo thang điểm $100$. Kết quả được trình bày trong bảng phân bố tần số ghép lớp sau đây:
	\begin{center}
		\begin{tabular}{|c|l|l|l|l|l|}
			\hline Số  điểm & {$[50; 60)$} & {$[60; 70)$} & {$[70; 80)$} & {$[80; 90)$} & {$[90; 100)$} \\
			 \hline Số người & $4$ & $7$ & $9$ & $18$ & $12$ \\
			\hline 
		\end{tabular}
	\end{center}
	Hãy ước lượng số trung  bình và mốt của mẫu số liệu ghép nhóm ở trên.
	 \shortans{2}
	 \loigiai{
		\begin{center}
			\begin{tabular}{|l|l|l|l|l|l|}
				\hline Số điềm & {$[50; 60)$} & {$[60; 70)$} & {$[70; 80)$} & {$[80; 90)$} &  {$[90; 100)$} \\
				\hline \begin{tabular}{l}
				Giá trị đại \\ 
				diện 
			\end{tabular} & $55$ & $65$ & $75$ & $85$ & $95$ \\
			\hline Số người & $4$ & $7$ & $9$ &  $18$ & $12$ \\
			 \hline
		\end{tabular}
	\end{center}
	\begin{itemize}
		\item  Điểm trung bình đánh giá phim Lật Mặt $6$ bằng:
		 $$
		\dfrac{55 \cdot 4+65 \cdot 7+75 \cdot  9+85 \cdot 18+95 \cdot 12}{50}=80{,}4
		 $$
		\item Nhóm chứa mốt của mẫu số liệu trên  là $[80; 90)$.\\
		Do  đó: $u_m=80; n_{m-1}=9; n_{m+1}=12; u_{m+1}-u_m=90-80=10$.
	\end{itemize}
	Mốt của  mẫu số liệu ghép nhóm là: $M_o=80+\dfrac{18-9}{(18-9)+(18-12)} \cdot 10=80{,}6$.
 }
\end{ex}

% nguồn: Ok_TraLoiNgan_Lop11_Lan1
\begin{ex}%[1D5H1-4]%[ma:65586616]
	Khảo sát số lần sử dụng Facebook của một người thực hiện mỗi ngày trong 30  ngày được lựa chọn ngẫu nhiên được thống kê trong bảng sau:\\
	\centerline{\begin{tabular}{|c|c|c|c|c|c|}
	\hline 
	Số lần sử dụng Facebook & $[3;5]$ & $[6;8]$ & $[9;11]$ & $[12;14]$  & $[15;17]$\\
	 \hline
	Số ngày & $2$ & $5$ & $11$ & $8$  & $4$\\
	 \hline
\end{tabular}}\\
Xác định mốt của mẫu số  liệu ghép nhóm trên.
\shortans{10,5} 
 \loigiai{
	Do số lần sử dụng Facebook là số nguyên nên  ta chỉnh lại như sau.\\
	\centerline{\begin{tabular}{|c|c|c|c|c|c|}
	 \hline
	Số lần sử dụng Facebook & $[2{,}5;5{,}5)$  & $[5{,}5;8{,}5)$ & $[8{,}5;11{,}5)$ & $[11{,}5;14{,}5)$ & $[14{,}5;17{,}5)$\\
	 \hline
	Số ngày & $2$ & $5$ & $11$ & $8$ & $4$\\ 
	\hline 
	\end{tabular}}\\
	Nhóm chứa  mốt của mẫu số liệu trên là $[8{,}5;11{,}5)$.\\
	Do đó $u_m=8{,}5$; $n_{m-1}=5$; $n_{m+1}=8$;  $u_{m+1}-u_m=11{,}5-8{,}5=3$.\\
	Mốt của mẫu số liệu ghép nhóm là $\mathrm{M_o}=8{,}5+\dfrac{11-5}{(11-5)+(11-8)}\cdot  3=10{,}5$\\
 }
\end{ex}

% nguồn: Ok_TraLoiNgan_Lop11_Lan1
\begin{ex}%[1D5H1-4]%[ma:af662d93]
	Kết quả kiểm tra môn Toán của lớp $11D$ như sau 
	\[ 
	\begin{array}{cccccccccccccccccccc}
		5 & 6 & 7 & 5 & 6 & 9 & 10 & 8 & 5 & 5 & 4 & 5 & 4 & 5 & 7 &  4 & 5 & 8 & 9 & 10 \\
		5 & 3 & 5 & 6 & 5 & 7 & 5 & 8 & 4 & 9 & 5 & 6 & 5 & 6 & 8 & 8 & 7 & 9 & 7 & 9 
	\end{array}
	 \]
	Mốt của bảng số liệu ghép nhóm trên là  bao nhiêu (làm tròn kết quả đến hàng phần mười)?
	 \shortans{6,2}
	\loigiai{\immini{Ta thấy: Nhóm $2$ ứng với nửa khoảng $\left[5;7\right)$ là nhóm có tần số lớn nhất với $u=5$, $g=2$, $n_2 = 18$. Nhóm $1$ có tần số $n_1  = 5$, nhóm $3$ có tần số $n_3=10$.\\
		Khi đó, mốt của mẫu số liệu là 
		 \[
		M_o = 5 + \left( \dfrac{18- 5}{2\cdot  18 - 5 - 10} \right) \cdot 2 \approx 6{,}2.
		 \]
		}{ 
		\begin{tabular}{|c|c|}
			 \hline
			\textbf{Nhóm} &  \textbf{Tần số}\\
			 \hline
			$\left[3;5\right)$ & $5$\\ 
			 \hline
			 $\left[5;7\right)$ & $18$\\
			\hline 
			 $\left[7;9\right)$ & $10$\\
			 \hline
			$\left[9;11\right)$ & $7$\\ 
			\hline 
			 & $n = 40$ \\
			 \hline
		\end{tabular}
		 }
	} 
\end{ex}

% nguồn: Ok_TraLoiNgan_Lop11_Lan1
\begin{ex}%[1D5H1-4]%[ma:3cd9a006]
	Người ta ghi lại tuổi thọ  của một số con ong cho kết quả như sau
	\begin{center}
		\begin{tabular}{|c|c|c|c|c|c|}
			 \hline
			Tuổi thọ (ngày)  & $[0;20)$ &$[20;40)$&$[40;60)$&$[60;80)$& $[80;100)$\\
			\hline 
			 Số lượng&$5$&$12$&$23$&$31$&$29$\\
			 \hline
		\end{tabular}
	\end{center}
	Tìm mốt của  mẫu số liệu.
	 \shortans{76}
	\loigiai{ 
		Tần số lớn nhất là  $31$ nên nhóm chứa mốt là nhóm $[60;80)$.\\
		 Ta có $j=4$; $a_4=60$; $m_4=31$; $m_3=23$; $m_5=29$; $h=20$.\\
		Khi  đó $M_O=60+\dfrac{31-23}{(31-23)+(31-29)}\cdot 20=76$.
	} 
\end{ex}

% nguồn: Ok_TraLoiNgan_Lop11_Lan1
\begin{ex}%[1D5H1-4]%[ma:479d35cf]
	Tuổi thọ (năm) của 50 bình ắc quy ô tô được cho  như sau
	\begin{center}
		\begin{tabular}{|c|c|c|c|c|c|c|}
			\hline 
			 Tuổi thọ (năm) & $\left[2;2{,}5 \right)$ & $\left[2{,}5;3 \right)$ & $\left[3;3{,}5 \right)$&$\left[3{,}5;4 \right)$&$\left[4;4{,}5 \right)$&$\left[4{,}5;5 \right)$ \\
			 \hline
			Tần số &$4$ & $9$ & $14$  &$11$ & $7$&$5$ \\
			 \hline
		\end{tabular}
	\end{center}
	Xác định mốt (làm tròn đến  hàng phần mười).
	\shortans{3,3} 
	\loigiai{ 
		Tần số lớn nhất  là $14$ nên nhóm chứa mốt là nhóm $\left[3;3, 5\right)$.\\
		Ta có $j=3$, $a_3=3$, $m_3=14$,  $m_2=9$, $m_4=11$, $h=0$, $5$. Do đó
		$$M_o=3+\dfrac{14-9}{(14-9)+(14-11)}\cdot 0{,}5\approx 3{,}3.$$ 
	 }
\end{ex}

% nguồn: Ok_TraLoiNgan_Lop11_Lan2
\begin{ex}%[1D5H1-4]%[ma:a4ad6c1b]
	Doanh thu bán hàng trong $20$ ngày được lựa chọn ngẫu nhiên của một của hàng được ghi lại ở bảng sau (đơn vị: triệu  đồng):
	\begin{center}
		\begin{tabular}{|c|c|c|c|c|c|}
			 \hline
			Doanh thu (triệu đồng) & [5; 7) & [7; 9) &  [9; 11) & [11; 13) & [13; 15) \\ \hline
			Số ngày & 2 & 7 &  7 & 3 & 1 \\ \hline
		\end{tabular}
	\end{center}
	Mốt của mẫu số liệu  trên là
	 \shortans{9}
	 \loigiai{
		Có $2$ nhóm chứa mốt của mẫu số liệu trên là  $[7; 9)$ và $[9; 11)$, do đó:\\
		Xét  nhóm $[7; 9)$ ta có:
		$M_o = 7 + \dfrac{7 - 2}{(7 - 2) +  (7 - 7)}\cdot(9 - 7) = 9$\\
		 Xét nhóm $[9; 11)$ ta có
		$M_o' = 9 + \dfrac{7 - 7}{(7  - 7) + (7 - 3)}\cdot(11 - 9) = 9$\\
		 Vậy mốt của mẫu số liệu là $9$.
	} 
\end{ex}

% nguồn: Ok_TraLoiNgan_Lop11_Lan2
\begin{ex}%[1D5H1-4]%[ma:52f9c7f4]
	Khảo sát số lần sử dụng Facebook của một người thực hiện mỗi ngày trong $30$  ngày được lựa chọn ngẫu nhiên được thống kê trong bảng sau
	\begin{center}
		\begin{tabular}{|l|c|c|c|c|c|c|c|}
			\hline Số lần sử dụng FB & {$[3;5])$} & {$ [6;8]$} & {$[9;11] $}  & {$[12;14]$} & {$[15;17]$}\\
			\hline Số ngày& $2$ & $5$ &  $11$ & $8$& $4$\\
			 \hline
		\end{tabular}
	\end{center}
	Tìm mốt của mẫu số  liệu ghép nhóm trên?
	\shortans{10,5} 
	 \loigiai{
		Do số  lần sử dụng Facebook là số nguyên nên ta chỉnh lại như sau
		\begin{center}
			\begin{tabular}{|l|c|c|c|c|c|c|c|}
				\hline Số lần sử dụng FB & {$[2{,}5;5{,}5]$} & {$ [5{,}5;8{,}5]$} & {$[8{,}5;11{,}5]  $} & {$[11{,}5;14{,}5]$} & {$[14{,}5;17{,}5]$}\\
				 \hline Số ngày& $2$ & $5$ & $11$ & $8$& $4$\\
				 \hline
			\end{tabular}
		\end{center}
		Nhóm chứa mốt  là ${[8{,}5; 11{,}5)}$ do đó ${u_m=8{,}5; n_{m-1}=5; n_{m+1}=8; u_{m+1}-u_m=11{,}5-8{,}5=3}$.\\
		mốt của mẫu số liệu ghép  nhóm là $$M_o=8{,}5+\dfrac{11-5}{(11-5)+(11-8)}\cdot 3=10{,}5.$$
	 }
\end{ex}

% nguồn: Ok_TraLoiNgan_Lop11_Lan2
\begin{ex}%[1D5H1-4]%[ma:955e7152]
	Một công ty bất động sản Đất Vàng thực hiện cuộc khảo sát khách hàng xem họ có nhu cầu mua nhà ở mức giá nào để tiến hành dự án  xây nhà ở Thăng Long group sắp tới. Kết quả khảo sát $500$ khách hàng được ghi lại ở bảng sau
	\begin{center}
		\begin{tabular}{|l|c|c|c|c|c|c|c|}
			\hline Mức giá (Triệu đồng) & {$[10;14)$}  & {$ [14;18)$} & {$[18;22) $} & {$[22;26)$} & {$[26;30)$}\\
			\hline Số  khách hàng & $75$ & $105$ & $179$ & $96$& $45$\\
			 \hline
		\end{tabular}
	\end{center}
	Công ty bất động sản Đất Vàng  nên xây nhà ở mức giá nào để nhiều người có nhu cầu xây nhà? (Kết quả làm tròn đến hàng phần chục)
	 \shortans{19,9}
	 \loigiai{
		Nhóm chứa  mốt của mẫu số liệu trên là ${[18; 22)}$.\\
		Do đó ${u_m=18; n_{m-1}=105;  n_{m+1}=96; u_{m+1}-u_m=22-18=4}$.\\
		mốt  của mẫu số liệu ghép nhóm là $$M_o=18+\dfrac{179-105}{(179-105)+(179-96)}\cdot 4\approx 19{,}9.$$\\
		Dựa vào kết quả trên ta có thể dự đoán rằng nếu công ty bất động sản Đất Vàng xây nhà ở mức giá $19{,}9$ triệu  đồng/mét vuông thì sẽ có nhiều người mua nhất.
	 }
\end{ex}

% nguồn: Ok_TraLoiNgan_Lop11_Lan2
\begin{ex}%[1D5H1-4]%[ma:45846b26]
	Khảo sát thời gian tập thể dục của  một số học sinh khối $11$ thu được mẫu số liệu ghép nhóm sau:
	\begin{center}
		\begin{tabular}{|c|c|c|c|c|c|}
			\hline 
			Thời gian (phút) & [0; 20) & [20; 40) & [40; 60) & [60; 80) & [80; 100) \\  \hline
			Số học sinh & 5 & 9 &  12 & 10 & 6 \\ \hline
		\end{tabular}
	\end{center}
	Mốt của  mẫu số liệu trên là
	 \shortans{52}
	\loigiai{ 
		Mốt $M_o$  chứa trong nhóm $[40; 60)$.\\
		Ta có : $ u_m  = 40;u_{m+1} = 60;u_{m+1}-u_m =60-40=20\\
		n_{m-1} =9; n_{m} =12; n_{m+1} = 10 $ \\ Do đó 
		$M_o = 40 + \dfrac{12 - 9}{(12 - 9) +  (12 - 10)}\cdot(60 - 40) = 52$
	 }
\end{ex}

% nguồn: Ok_TraLoiNgan_Lop11_Lan2
\begin{ex}%[1D5H1-4]%[ma:f2d6a3af]
	Khảo sát thời gian tự học bài ở nhà của một số em học sinh lớp $11$ thu  được mẫu ghép nhóm số liệu như sau:
	\begin{center}
		\begin{tabular}{|c|c|c|c|c|c|}
			\hline 
			Thời gian (phút) & [0; 30)  & [30; 60) & [60; 90) & [90; 120) & [120; 150) \\ \hline
			Số học sinh & 9 & 10  & 9 & 15 & 7 \\ \hline
		\end{tabular}
	\end{center}
	Tìm mốt của số  liệu ghép nhóm trên (làm tròn kết quả đến hàng đơn vị)
	\shortans{102} 
	 \loigiai{
		Nhóm chứa  mốt của mẫu số liệu là nhóm $[90; 120)$. \\Vậy mốt
		$M_o = 90 + \dfrac{15 - 9}{15 - 9 + 15 - 7} \cdot  30 = 102{,}86\approx 102$.
	} 
\end{ex}

% nguồn: Ok_TraLoiNgan_Lop11_Lan2
\begin{ex}%[1D5H1-4]%[ma:f2eb9467]
	Người ta tiến hành phỏng vấn $50$ người về phim chiếu rạp Lật mặt $6$ của Lý Hải. Người điều tra yêu cầu cho điểm phim theo  thang điểm $100$. Kết quả được trình bày trong bảng phân bố tần số ghép lớp sau đây
	\begin{center}
		\begin{tabular}{|l|c|c|c|c|c|c|c|}
			\hline Số điểm & {$[50;60)$} & {$ [60;70)$} & {$[70;80) $}  & {$[80;90)$} & {$[90;100)$}\\
			\hline Số người & $4$ & $7$ &  $9$ & $18$& $12$\\
			 \hline
		\end{tabular}
	\end{center}
	Hãy  tìm mốt của mẫu số liệu ghép nhóm ở trên? (Kết quả làm tròn đến hàng phần chục)
	 \shortans{80,6}
	 \loigiai{
		Nhóm chứa mốt của mẫu số liệu  trên là $[80;90)$.\\
		Do đó  ${u_m=80; n_{m-1}=9; n_{m+1}=12; u_{m+1}-u_m=90-80=10}$.\\
		mốt của mẫu số  liệu ghép nhóm là $M_o=80+\dfrac{18-9}{(18-9)+(18-12)}\cdot 10=80{,}6$.
	} 
\end{ex}

% nguồn: Ok_TraLoiNgan_Lop11_Lan2
\begin{ex}%[1D5H1-4]%[ma:ddc7398d]
	Kết quả điểm thi của $32$ học sinh trong kì thi Tiếng Anh  (thang điểm $100$) được biểu diễn ở biểu đồ dưới đây
	\begin{center}
		\begin{tikzpicture}[scale=0.8,>=stealth,line join=round,line cap=round,font=\footnotesize,scale=0.85,line width=1pt]
				\draw[->]  (0,0)--(0,9);
				\foreach \x/\luongmua  in {0/0\%,1/5\%,2/10\%,3/15\%,4/20\%,5/25\%,6/30\%,7/35\%}{
					\draw (0,\x)  node[label=left:\textbf{{\luongmua}}] {};
				} 
				\foreach  \y in {1,2,3,4,5,6,7}{
					\draw[dashed,thin,line  width=0.01pt] (0,\y)--(13.5,\y);
				 }
				\draw[line cap=round,pattern=north east lines] (1,0)--(1,2.5)--(2,2.5) node[midway,above]{$13\%  $}--(2,0)--cycle;
				\draw[line  cap=round,pattern=north east lines] (3,0)--(3,3.8)--(4,3.8) node[midway,above]{$19\% $}--(4,0)--cycle;
				\draw[line cap=round,pattern=north east lines] (5,0)--(5,6.2)--(6,6.2)  node[midway,above]{$31\% $}--(6,0)--cycle;
				\draw[line cap=round,pattern=north  east lines] (7,0)--(7,3.8)--(8,3.8) node[midway,above]{$19\% $}--(8,0)--cycle;
				\draw[line cap=round,pattern=north  east lines] (9,0)--(9,2.5)--(10,2.5) node[midway,above]{$13\% $}--(10,0)--cycle;
				\draw[line cap=round,pattern=north east lines]  (11,0)--(11,1.2)--(12,1.2) node[midway,above]{$6\% $}--(12,0)--cycle;
				\node [below]  at (1.5,0){$[40; 50)$};
				\node [below]  at (3.5,0){$[50; 60)$};
				\node [below] at (5.5,0){$[60; 70)$}; 
				\node [below]  at (7.5,0){$[70; 80)$};
				\node [below] at  (9.5,0){$[80; 90)$};
				\node [below]  at (11.5,0){$[90; 100)$};
				\draw[->] (0,0) node [below]{$ O $}--(13.5,0)  node[above]{$x$(gam)};
			\end{tikzpicture}
	\end{center}
	 Hãy tìm mốt của mẫu số liệu trên?
	\shortans{65} 
	\loigiai{ 
		Ta có bảng tần số ghép  nhóm sau
		\begin{center}
			\begin{tabular}{|l|c|c|c|c|c|c|c|}
				\hline Lớp điểm & {$[40;50)$}& {$[50;60)$} & {$  [60;70)$} & {$[70;80) $} & {$[80;90)$} & {$[90;100)$}\\
				\hline Số người & $4$ & $6$ & $10$ & $6$& $4$& $2$\\ 
				\hline 
			\end{tabular}
		\end{center}
		Nhóm chứa mốt  của mẫu số liệu là $[60;70)$.\\
		Do đó mốt của mẫu số liệu là $60+\dfrac{10-6}{10-6+10-6}\cdot  10=65$
	} 
\end{ex}

% nguồn: Ok_TraLoiNgan_Lop11_Lan2
\begin{ex}%[1D5H1-4]%[ma:0c53a459]
	Điểm kiểm tra giữa kì I của một lớp được cô giáo chủ nhiệm ghi  lại theo bảng sau:
	\begin{center}
		\begin{tabular}{|c|c|c|c|c|c|}
			\hline 
			Điểm & [4; 5,5) & [5,5; 6,5) & [6,5; 7,5) &  [7,5; 8,5) & [8,5; 9,5) \\ \hline
			Số học sinh & 3 & 7 & 9 & 15 &  6 \\ \hline
		\end{tabular}
	\end{center}
	Mốt của  mẫu số liệu ghép nhóm trên là
	\shortans{7,9} 
	 \loigiai{
		Nhóm chứa mốt của mẫu số liệu là  nhóm $[7{,}5; 8{,}5)$. \\Vậy mốt $M_{o} =7{,}5 + \dfrac{15 - 9}{15 - 9 + 15 - 6} \cdot 1 = 7{,}9$.
	 }
\end{ex}

% nguồn: Ok_TraLoiNgan_Lop11_Lan2
\begin{ex}%[1D5H1-4]%[ma:64bbd198]
	Kết quả điểm thi của $32$ học sinh trong kì thi Tiếng Anh (thang  điểm $100$) được biểu diễn ở biểu đồ dưới đây:
	\begin{center}
		\begin{tikzpicture}[>=stealth,declare function={ox=-0.5; kc=1.3; rc=0.6; kcy=0.9;}]
				\foreach[count=\i] \yt in  {5,10,15,20,25,30,35}{
					\draw[gray!35]  (0,{\i*kcy}) node[left,text=black]{$\yt\,\%$}--+(9,0);
				} 
				\foreach[count=\i] \yt/\xt/\gt in {2.6/[40; 50)/13\%,3.8/[50; 60)/19\%,6.2/[60;  70)/31\%,3.8/[70; 80)/19\%,2.6/[80; 90)/13\%,1.2/[90; 100)/6\%}{
					\draw[blue,fill=blue!25] ({ox+\i*kc},0)  rectangle +(rc,{\yt*kcy});
					 \path ({ox+\i*kc+0.5*rc},0) node[below]{\xt};
					 \path ({ox+\i*kc+0.5*rc},{\yt*kcy}) node[above]{$\gt$};
				} 
				\draw[->] (0,0)--(8.5,0)  node[midway,shift={(0,-25pt)}]{Điểm};
				\draw[->] (0,0)--(0,6.75)  node[left,anchor=east]{};
			\end{tikzpicture}
	\end{center}
	Mốt của  mẫu số liệu trên là
	\shortans{65} 
	 \loigiai{
		Ta có  bảng tần số ghép nhóm sau:
		\begin{center}
			\begin{tabular}{|c|c|c|c|c|c|c|}
				 \hline
				Lớp điểm & [40; 50) & [50; 60) & [60; 70) & [70; 80) & [80;  90) & [90; 100) \\ \hline
				Tần số & 4 & 6 & 10  & 6 & 4 & 2 \\ \hline
			\end{tabular}
		\end{center}
		Nhóm chứa mốt của mẫu số liệu  là nhóm $[60; 70)$.\\ Vậy mốt $M_{o}= 60 + \dfrac{10 - 6}{10 - 6 + 10 - 6} \cdot 10 = 65$
	} 
\end{ex}

% nguồn: Ok_TraLoiNgan_Lop11_Lan2
\begin{ex}%[1D5H1-4]%[ma:33f53d70]
	Số sản phẩm bán được của một  hãng X tại một cửa hàng trong $30$ ngày được thống kê trong bảng sau:
	\begin{center}
		\begin{tabular}{|c|c|c|c|c|c|c|}
			 \hline
			Số sản phẩm  & [3; 5] & [6; 8] & [9; 11] & [12; 14] & [15; 17] \\ \hline
			Số ngày & 3 & 5 & 7 & 10 & 5 \\  \hline
		\end{tabular}
	\end{center}
	Cửa hàng đó bán được cao nhất bao nhiêu sản phẩm mỗi ngày? (kết  quả làm tròn đến hàng phần chục)
	\shortans{12,8} 
	 \loigiai{
		Nhóm chứa  mốt của mẫu số liệu là $[12; 14]$.
		Mốt của mẫu số liệu là:\\ $M_{o}=12 + \dfrac{10 - 7}{10 - 7 +  10 - 5} \cdot 2 = 12{,}75\approx 12{,}8$.\\
		Vậy cửa hàng đó bán được cao nhất  $12$ sản phẩm mỗi ngày.
	 }
\end{ex}

% nguồn: Ok_TraLoiNgan_Lop11_Lan2
\begin{ex}%[1D5H1-4]%[ma:f375fce2]
	Kiểm tra cân nặng $25$ quả dưa của hai nhà vườn A  và B cho kết quả như bảng sau:
	\begin{center}
		\begin{tabular}{|c|c|c|c|c|c|}
			\hline 
			Cân  nặng (kg) & [1,5; 2,5) & [2,5; 3,5) & [3,5; 4,5) & [4,5; 5,5) & [5,5; 6,5) \\ \hline
			Số quả của vườn A & 2 & 6 & 4 & 10 & 3 \\ \hline 
			 Số quả của vườn B & 3 & 7 & 8 & 5 & 2 \\ \hline
		\end{tabular}
	\end{center}
	Nếu lấy một quả  dưa bất kì của vườn A thì khả năng cao nhất là lấy được quả dưa nặng bao nhiêu kg?
	\shortans{4,96} 
	 \loigiai{
		Mốt cân nặng của các quả dưa vườn A thuộc nhóm $[4{,}5; 5{,}5)$. Vậy mốt cân nặng của các quả dưa vườn A  là
		$M_o = 4{,}5 +  \dfrac{10 - 4}{10 - 4 + 10 - 3} \cdot 1 = 4{,}96$.\\
		Dựa vào kết quả trên, khả năng lấy được quả dưa nặng $4{,}96 $ kg  là cao nhất.
	 }
\end{ex}

% nguồn: Ok_TraLoiNgan_Lop11_Lan2
\begin{ex}%[1D5H1-4]%[ma:128c2080]
	Một đại lí bảo hiểm đã thống kê số lượng khách mua bảo hiểm nhân thọ trong một ngày ở biểu đồ sau 
	\begin{center}
		\begin{tikzpicture}[scale=0.8,>=stealth,line join=round,line cap=round,font=\footnotesize,scale=0.85,line width=1pt]
				\draw[->] (0,0)--(0,11)  node[left]{\text{Người}};
				\foreach \y in {1,2,3,4,5,6,7,8,9,10} 
				\draw[shift={(0,\y)}] (0,0)--(-2pt,0)  node[left]{\scriptsize ${\y}$};
				\path  (6,-1.5) node {
					$\begin{array}{c}
					\normalsize{\textbf{Buổi đồ số lượng khách hàng theo giới tính  và độ tuổi}}
				\end{array}$
					}; 
				\foreach \y in  {1,2,3,4,5,6,7,8,9,10}{
					\draw[dashed,thin,line  width=0.01pt] (0,\y)--(13,\y);
				} 
				\draw[line cap=round,pattern=north east lines]  (10,6)--(10,7)--(9,7)--(9,6)--(10,6) node[above right]{\text{Nam}};
				 \draw[line cap=round,pattern=dots] (10,8)--(10,9)--(9,9)--(9,8)--(10,8) node[above right]{\text{Nữ}};
				\draw[line cap=round,pattern=north east lines]  (0,0)--(0,4)--(1,4) node[midway,above]{$ $}--(1,0)--cycle;
				 \draw[line cap=round,pattern=dots] (1,0)--(1,3)--(2,3) node[midway,above]{$ $}--(2,0)--cycle;
				 \draw[line cap=round,pattern=north east lines] (2,0)--(2,6)--(3,6) node[midway,above]{$ $}--(3,0)--cycle;
				\draw[line cap=round,pattern=dots]  (3,0)--(3,8.7)--(4,8.7) node[midway,above]{$ $}--(4,0)--cycle;
				 \draw[line cap=round,pattern=north east lines] (4,0)--(4,10)--(5,10) node[midway,above]{$ $}--(5,0)--cycle;
				\draw[line cap=round,pattern=dots] (5,0)--(5,6)--(6,6) node[midway,above]{$  $}--(6,0)--cycle;
				\draw[line cap=round,pattern=north east  lines] (6,0)--(6,7.5)--(7,7.5) node[midway,above]{$ $}--(7,0)--cycle;
				\draw[line cap=round,pattern=dots] (7,0)--(7,4)--(8,4) node[midway,above]{$  $}--(8,0)--cycle;
				\draw[line cap=round,pattern=north east lines] (8,0)--(8,3)--(9,3) node[midway,above]{$ $}--(9,0)--cycle; 
				\draw[line cap=round,pattern=dots] (9,0)--(9,2)--(10,2) node[midway,above]{$  $}--(10,0)--cycle;
				\node [below] at  (1,0){$ [20; 30)$};
				\node [below] at (3,0){$  [30; 40)$};
				\node [below] at (5,0){$  [40; 50)$};
				\node [below] at (7,0){$ [50; 60)$}; 
				 \node [below] at (9,0){$ [60; 70)$};
				\draw[->] (0,0) node  [below left=-2pt]{$ O $}--(13,0) node[below]{\text{Tuổi}};
			\end{tikzpicture}
	\end{center}
	Hãy sử dụng dữ liệu trên để tư vấn cho đại lí bảo hiểm xác định khách hàng nam và nữ ở độ tuổi nào hay mua bảo hiểm nhất (kết quả  làm tròn đến hàng đơn vị)
	 \shortans{37}
	 \loigiai{
		mốt của mẫu số  liệu đối với khách hàng nam là $M_o=40+\dfrac{10-6}{(10-6)+(10-7)} \cdot 10\approx 46$.\\
		Nên với nam ở độ tuổi từ $40$ đến $50$ có nhu cầu mua bảo hiểm  lớn nhất, đặc biệt là độ tuổi $46$.\\
		mốt của mẫu số  liệu đối với khách hàng nữ là $M_o=30+\dfrac{9-3}{(9-3)+(9-6)}\cdot10\approx 37$.\\
		Nên với nử ở độ tuổi từ $30$ đến  $40$ có nhu cầu mua bảo hiểm lớn nhất, đặc biệt là độ tuổi $37$.
	 }
\end{ex}

% nguồn: Ok_TraLoiNgan_Lop11_Lan2
\begin{ex}%[1D5H1-4]%[ma:52c334c7]
	Cho mẫu số liệu ghép nhóm về chiều cao của $25$ cây dừa giống  như sau:
	\begin{center}
		\begin{tabular}{|c|c|c|c|c|c|}
			 \hline
			Chiều cao (cm) & [0; 10) & [10; 20) & [20; 30) & [30; 40) & [40; 50) \\ \hline 
			Số cây & 4 & 6 & 7 & 5 & 3  \\ \hline
		\end{tabular}
	\end{center}
	Mốt của mẫu số liệu  ghép nhóm này là (làm tròn kết quả đến hàng phần chục)
	 \shortans{23,3}
	\loigiai{ 
		Tần số lớn nhất là $7$, nên nhóm chứa  mốt là nhóm $[20; 30)$. \\Vậy
		$M_o = 20 + \dfrac{7 - 6}{(7 - 6) + (7 - 5)}  \cdot 10 = 23{,}3$.
	} 
\end{ex}

% nguồn: Ok_TraLoiNgan_Lop11_Lan2
\begin{ex}%[1D5H1-4]%[ma:a0195e73]
	Tuổi thọ (năm) của $50$ bình ắc quy ô tô  được thu thập theo mẫu số liệu sau:
	\begin{center}
		\begin{tabular}{|c|c|c|c|c|c|}
			\hline 
			Tuổi thọ (năm) & [2; 2,5) & [2,5;  3) & [3; 3,5) & [4; 4,5) & [4,5; 5,5) \\ \hline
			Tần số & 4 & 9 & 14 & 11 &  7 \\ \hline
		\end{tabular}
	\end{center}
	Xác định mốt? (kết quả làm tròn đến  hàng phần trăm)
	\shortans{3,31} 
	 \loigiai{
		Nhóm chứa mốt là nhóm  $[3; 3{,}5)$. Do đó
		$M_o = 3 + \dfrac{14 -  9}{14 - 9 + 14 - 11} \cdot 0{,}5 \approx 3{,}31$
	} 
\end{ex}

% nguồn: Ok_TraLoiNgan_Lop11_Lan2
\begin{ex}%[1D5H1-4]%[ma:e309a3bc]
	 Kết quả đo chiều cao của $250$ cây dừa đột biến $3$ năm tuổi ở một viện nghiên cứu được tổng hợp ở bảng sau
	\begin{center}
		\begin{tabular}{|l|c|c|c|c|c|c|c|}
			\hline Chiều  cao $\left(m^2\right) $ & {$[8{,}5;8{,}8)$} & {$ [8{,}8;9{,}1)$} & {$[9{,}1;9{,}4) $} & {$[9{,}4;9{,}7)$} & {$[9{,}7;10)$}\\
			\hline Số cây & $36$ & $45$ & $83$ & $65$& $21$\\ 
			 \hline
		\end{tabular}
	\end{center}
	Hãy tìm mốt của mẫu số liệu ghép nhóm trên? (kết quả làm  tròn đến hàng phần chục)
	 \shortans{9,3}
	 \loigiai{
		Nhóm chứa mốt của mẫu số  liệu trên là ${[9{,}1; 9{,}4)}$\\
		Do đó: ${u_m=9{,}1; n_{m-1}=45; n_{m+1}=65;  u_{m+1}-u_m=9{,}4-9{,}1=0{,}3}$.\\
		mốt của  mẫu số liệu ghép nhóm là $M_o=9{,}1+\dfrac{83-45}{(83-45)+(83-65)}\cdot 0{,}3\approx 9{,}3.$
	 }
\end{ex}

% nguồn: Ok_TraLoiNgan_Lop11_Lan1
\begin{ex}%[1D5V1-4]%[ma:d5b8eb46]
	 Một công ty bất động sản Đất Vàng thực hiện cuộc khảo sát khách hàng xem họ có nhu cầu mua nhà ở mức giá nào để tiến hành dự án xây nhà ở Thăng Long group sắp tới. Kết quả khảo sát $500$ khách hàng được ghi lại ở bảng sau:
	\begin{center}
		\begin{tabular}{|l|c|c|c|c|c|}
			\hline Mức giá (triệu đồng) & {$[10; 14)$} & {$[14; 18)$} & {$[18; 22)$} & {$[22; 26)$}  & {$[26; 30)$} \\
			\hline  Số khách hàng & $75$ & $105$ & $179$ & $96$ & $45$ \\
			\hline 
		\end{tabular}
	\end{center}
	Công ty bất động sản Đất Vàng nên xây  nhà ở mức giá nào để nhiều người có nhu cầu xây nhà?
	 \shortans{19,9}
	 \loigiai{
		Nhóm chứa mốt của mẫu số liệu  trên là $[18; 22)$.\\
		 Do đó $u_m=18$; $n_{m-1}=105$; $n_{m+1}=96$; $u_{m+1}-u_m=22-18=4$.\\
		 Mốt của mẫu số liệu ghép nhóm là $M_0=18+\dfrac{179-105}{(179-105)+(179-96)} \cdot 4 \approx 19{,}9$.\\
		Dựa vào kết quả trên ta có thể dự đoán rằng nếu công ty bất động sản Đất Vàng xây nhà ở mức giá $19{,}9$ triệu đồng$/\mathrm{m}^2$ thì  sẽ có nhiều người mua nhất.
	} 
\end{ex}

% nguồn: Ok_TraLoiNgan_Lop11_Lan1
\begin{ex}%[1D5V1-4]%[ma:f9db41a8]
	Một công ty xây dựng khảo sát khách hàng xem họ có nhu cầu mua nhà ở mức giá nào. Kết quả được khảo sát được ghi lại ở  bảng sau
	\begin{center}
		\begin{tabular}{|c|c|c|c|c|c|}
			\hline 
			Mức giá (triệu đồng/$m^2$) & $\left[10;14\right)$ & $\left[14;18\right)$ & $\left[18;22\right)$ &  $\left[22;26\right)$ & $\left[26;30\right)$ \\
			\hline 
			Số khách hàng & 54 & 78 & 120 & 45 &  12 \\
			 \hline
		\end{tabular}
	\end{center}
	Tìm mốt của mẫu số liệu ghép nhóm trên  (làm tròn kết quả đến hàng phần chục).
	\shortans{19,4} 
	 \loigiai{
		Nhắc lại  về cách tìm mốt của mẫu số liệu:
		\begin{enumerate}
			 \item Xác định nhóm có tần số lớn nhất (gọi là nhóm chứa mốt), giả sử là nhóm $j$: $\left[a_{j};a_{j+1}\right)$.
			\item Mốt được xác định là $M_0 = a_j + \dfrac{m_j - m_{j-1}}{\left(m_j - m_{j-1}\right)+\left(m_j - m_{j+1}\right)} \cdot h$. trong đó $m_j$ là tần số của nhóm $j$ (quy ước $m_0 = m_{k+1} = 0$) và $h$ là độ dài của  nhóm.
		\end{enumerate}
		Nhóm chứa mốt của số liệu trên là  nhóm $\left[18;22\right)$.\\
		Do đó $j = 3$, $a_3  = 18$, $m_3 = 120$, $m_2 = 78$, $m_4 = 45$, $h=4$.\\
		 Vậy mốt của mẫu số liệu ghép nhóm là
		\[M_0 = 18 + \dfrac{120-78}{\left(120-78\right)+\left(120-45\right)} \cdot 4 = \dfrac{758}{39} \approx 19{,}4.\] 
	} 
\end{ex}

% nguồn: Ok_TraLoiNgan_Lop11_Lan1
\begin{ex}%[1D5V1-4]%[ma:06906041]
	Hãy sử dụng dữ liệu ở để tư vấn cho đại lí bảo hiểm xác định khách  hàng nam ở tuổi nào hay mua bảo hiểm nhất. Số khách hàng mua bảo hiểm ở từng độ tuổi được thống kê như sau:
	\begin{center}
		\begin{tabular}{|l|c|c|c|c|c|}
			 \hline
			Độ tuổi & $[20; 30)$ & $[30; 40)$ & $[40; 50)$ & $[50; 60)$ &  $[60; 70)$ \\
			\hline 
			Số khách hàng nam & $4$ & $6$ & $10$  & $7$ & $3$ \\
			\hline 
		\end{tabular}
	\end{center}
	 \shortans{46}
	\loigiai{ 
		 Nhóm chứa mốt của mẫu số liệu khách hàng nam là $[40; 50)$.\\
		Do đó $u_m=40, n_{m-1}=6; n_{m+1}=7; u_{m+1}-u_m=50-40=10$.\\ 
		Mốt của mẫu số liệu  nhóm khách hàng nam là
		 $M_0=40+\dfrac{10-6}{(10-6)+(10-7)} \cdot 10=45{,}7$.\\
		Dựa vào kết quả trên ta có thể dự đoán được khách hàng nam $46$ tuổi có  nhu cầu mua bảo hiểm cao nhất.
	 }
\end{ex}


%%%%%%%%%%%%%%%% TỰ LUẬN %%%%%%%%%%%%%%%%

% nguồn: Ok_TuanLuan_Toan11_Lan1
\begin{bt}%[1D5N1-4]%[ma:7d8b07f8]
	Bảng số liệu ghép nhóm sau cho biết chiều cao (cm) của $50$ bạn học sinh lớp $11A$. 
	\begin{center}
		\begin{tabular}{|c|c|c|c|c|c|}
			\hline 
			Khoảng chiều cao (cm)& $[145;150)$  & $[150;155)$ & $[155;160)$ & $[160;165)$ & $[165;170)$ \\
			 \hline
			Số học sinh& $7$  & $14$ & $10$ & $10$ & $9$ \\
			 \hline
		\end{tabular}
	\end{center}
	Tính mốt của mẫu số  liệu ghép nhóm này. Có thể kết luận gì từ giá trị nhận được?
	\loigiai{ 
		Tần số lớn nhất là $14$ nên nhóm $[150;155)$  là nhóm chứa mốt.\\
		Ta có $j=2$; $a_2=150$; $m_2=14$; $m_1=7$; $m_3=10$; $h=5$.  Do đó
		 $$M_o=150+\dfrac{14-7}{(14-7)+(14-10)}\cdot5\approx153{,}18.$$
		Số học sinh có  chiều cao khoảng $153{,}18$ cm là nhiều nhất
	 }
\end{bt}

% nguồn: Ok_TuanLuan_Toan11_Lan2
\begin{ex}%[1D5N1-4]%[ma:72ac0aa8]
	Bảng số liệu ghép nhóm sau cho biết chiều cao (cm) của  $50$ bạn học sinh lớp $11A$.
	\begin{center}
		\begin{tabular}{|c|c|c|c|c|c|}
			 \hline
			Khoảng chiều cao (cm)& $[145;150)$ & $[150;155)$  & $[155;160)$ & $[160;165)$ & $[165;170)$ \\
			\hline 
			Số học sinh& $7$ & $14$ & $10$ & $10$  & $9$ \\
			\hline 
		\end{tabular}
	\end{center}
	Tính mốt của mẫu số liệu ghép nhóm này. Có  thể kết luận gì từ giá trị nhận được?
	 \loigiai{
		Tần số lớn nhất là  $14$ nên nhóm $[150;155)$ là nhóm chứa mốt.\\
		Ta có $j=2$; $a_2=150$; $m_2=14$; $m_1=7$; $m_3=10$;  $h=5$. Do đó
		 $$M_o=150+\dfrac{14-7}{(14-7)+(14-10)}\cdot5\approx153{,}18.$$
		Số học sinh có chiều cao khoảng  $153{,}18$ cm là nhiều nhất
	} 
\end{ex}

% nguồn: Ok_TuanLuan_Toan11_Lan2
\begin{ex}%[1D5N1-4]%[ma:dfbe394f]
	Điều tra  về thu nhập của công nhân xí nghiệp X (đơn vị: nghìn đồng/ tháng), người ta ghi được bảng tần số ghép lớp sau đây
	\begin{center}
		\begin{tabular}{|l|c|c}
			\hline  Tiền lương & Số người\\
			\hline $[800;890]$ & $15$  \\
			\hline  $[900;990]$ & $25$ \\
			\hline  $[1000;1090]$ & $28$ \\
			\hline $[1100;1190]$ &  $35$ \\
			\hline  $[1200;1290]$ & $40$ \\
			\hline  $[1300;1390]$ & $30$ \\
			\hline $[1400;1490]$  & $27$ \\
			\hline 
		\end{tabular}
	\end{center}
	Hãy tìm thu nhập trung bình và mốt của  mẫu số liệu ghép nhóm trên (kết quả làm tròn đến hàng phần trăm).
	 \loigiai{
		Hiệu  chỉnh lại bảng số liệu ta có
		\begin{center}
			\begin{tabular}{|l|c|c|c|c|c|c|c|}
				\hline  Tiền lương & {$[795;895)$} & {$ [895;995)$} & {$[995;1095) $} & {$[1095;1195)$}& {$[1195;1295)$}& {$[1295;1395)$}& {$[1395;1495)$}\\
				\hline GT đại diện & $845$ & $945$ & $1045$  & $1145$& $1245$& $1345$& $1445$\\
				\hline Số người &  $15$ & $25$ & $28$ & $35$& $40$& $30$& $27$\\
				 \hline
			\end{tabular}
		\end{center}
		Thu nhập trung bình của công  nhân bằng
		$$\overline{x}=\dfrac{845\cdot 15+945\cdot 25+1045\cdot 28+1145\cdot 35+1245\cdot 40+1345\cdot 30+1445\cdot 27}{200}=1174.$$ 
		Nhóm chứa mốt của  mẫu số liệu là $[1195;1295)$.\\
		Do đó $u_m=1195, n_{m-1}=35, n_{m+1}=30,  u_{m+1}-u_m=1295-1195=100$.\\
		mốt của mẫu số  liệu ghép nhóm là
		$$ M_o=1195+\dfrac{40-35}{(40-35)+(40-30)}\cdot 100=\dfrac{3685}{3}\approx  1228{,}33.$$
	 }
\end{ex}

% nguồn: Ok_TuanLuan_Toan11_Lan2
\begin{ex}%[1D5N1-4]%[ma:061a2ef7]
	Cân lần lượt $40$ quả cam (đơn vị gram) ta được kết  quả sau (mẫu số liệu)
	\begin{center}
		\begin{tabular}{|l|c|c|c|c|c|c|c|c|c|c|c|c|c||c|c|}
			\hline $85$ & $86$& $86$& $86$& $86$& $86$& $87$& $87$& $87$& $87$&  $87$& $88$& $88$& $88$& $88$& $89$\\
			\hline $89$  & $89$ &$89$ &$89$ &$89$ &$89$ &$90$ &$90$ &$90$ &$90$ &$90$ &$91$ &$91$ &$91$ &$92$ &$93$ \\
			\hline $93$ & $93$ &$93$ &$94$ & $94$ &$94$  &$94$ &$94$ &$94$ &$ $& $ $& $ $& $ $& $ $& $ $& $ $\\
			 \hline
		\end{tabular}
	\end{center}
	\begin{enumEX}{1}
		 \item Tính khối lượng trung bình của $40$ quả cam.
		\item Tổng hợp  lại số liệu trên vào bảng tần số ghép nhóm theo mẫu sau
		\begin{center}
			\begin{tabular}{|l|c|c|c|c|c|c|}
				\hline Khối  lượng & {$[85;87)$} & {$ [87;89)$} & {$[89;91) $} & {$[91;93)$}& {$[93;95)$}\\
				\hline Số quả & ? & ? & ?  & ? & ?\\
				 \hline
			\end{tabular}
		\end{center}
		\item Hãy tính khối lượng trung bình của $40$ quả cam từ  bảng tần số ghép nhóm trên.
		\item Dự  đoán khối lượng quả cam xuất hiện nhiều nhất (kết quả làm tròn đến hàng phần trăm).
	\end{enumEX}
	\loigiai{ 
		\begin{enumEX}{1}
			\item Khối lượng trung bình của 40 quả cam là $$\overline{x}=\dfrac{85+86\cdot 5+87\cdot 5+88\cdot 4+89\cdot 7+90\cdot 5+91\cdot  3+92+93\cdot 4+94\cdot 6}{40}=91{,}9\text{(g)}.$$
			\item Bảng tần  số ghép nhóm
			\begin{center}
				\begin{tabular}{|l|c|c|c|c|c|c|}
					\hline Khối lượng  & {$[85;87)$} & {$ [87;89)$} & {$[89;91) $} & {$[91;93)$}& {$[93;95)$}\\
					\hline Số quả & $6$ & $9$ & $11$ &  $4$& $10$\\
					 \hline
				\end{tabular}
			\end{center}
			\item Ta có bảng thống kê theo giá trị đại  diện
			\begin{center}
				\begin{tabular}{|l|c|c|c|c|c|c|}
					\hline Khối lượng & {$[85;87)$} & {$ [87;89)$} & {$[89;91) $} & {$[91;93)$}&  {$[93;95)$}\\
					\hline Giá trị đại  diện & $86$ & $88$ & $90$ & $92$& $94$\\
					\hline Số quả & $6$ & $9$  & $11$ & $4$& $10$\\
					 \hline
				\end{tabular}
			\end{center}
			Khối lượng trung bình của $40$ quả cam  là $$\overline{x}=\dfrac{86\cdot 6+88\cdot 9+90\cdot 11+92\cdot 4+94\cdot 10}{40}=90{,}15\text{(g)}.$$
			 \item Nhóm chứa mốt của mẫu số liệu trên là $[89;91)$.\\
			Do đó:  $u_m=89;n_{m-1}=9,n_{m+1}=4,u_{m+1}-u_m=91-89=2$.\\
			mốt của mẫu số  liệu ghép nhóm là $$ M_o=89+\dfrac{11-9}{(11-9)+(11-4)}\cdot 2=\dfrac{805}{9}\approx 89{,}44\text{(g)}.$$
			Vậy khả năng quả cam nặng $89{,}44\text{ (g)}$ là cao  nhất.
		\end{enumEX}
	 }
\end{ex}

% nguồn: Ok_TuanLuan_Toan11_Lan2
\begin{ex}%[1D5H1-4]%[ma:3426ebf9]
	Kết quả đo khối lượng của $30$ củ khoai tây ở nông trường được biểu diễn ở  biểu đồ dưới đây
	\begin{center}
		\begin{tikzpicture}[scale=0.8,>=stealth,line join=round,line cap=round,font=\footnotesize,scale=0.85,line width=1pt]
				\draw[->]  (0,0)--(0,9);
				\foreach  \x/\luongmua in {2/3,4/6,3.3/5,6/9,8/12}{
					 \draw (0,\x) node[label=left:\textbf{{\luongmua}}] {};
				 }
				\path (6,-1.5)  node {
					$\begin{array}{c}
					\normalsize{\textbf{Khối lượng $30$ củ khoai  tây}}
				\end{array}$
					}; 
				\foreach  \y in {2,3.3,4,8}{
					\draw[dashed,line width=0.01pt]  (0,\y)--(11,\y);
				 }
				\draw[line cap=round,pattern=north east lines] (1,0)--(1,2)--(2.5,2) node[midway,above]{$  $}--(2.5,0)--cycle;
				\draw[line  cap=round,pattern=north east lines] (2.5,0)--(2.5,3.3)--(4,3.3) node[midway,above]{$ $}--(4,0)--cycle;
				\draw[line cap=round,pattern=north east  lines] (4,0)--(4,8)--(5.5,8) node[midway,above]{$ $}--(5.5,0)--cycle;
				\draw[line cap=round,pattern=north east lines] (5.5,0)--(5.5,4)--(7,4)  node[midway,above]{$ $}--(7,0)--cycle;
				 \draw[line cap=round,pattern=north east lines] (7,0)--(7,2)--(8.5,2) node[midway,above]{$ $}--(8.5,0)--cycle;
				\node [below] at (1,0){$70$}; 
				 \node [below] at (2.5,0){$80$};
				 \node [below] at (4,0){$90$};
				\node  [below] at (5.5,0){$100$};
				\node  [below] at (7,0){$110$};
				\node [below]  at (8.5,0){$120$};
				\draw[->] (0,0) node [below left=-2pt]{$ O $}--(11,0) node[below]{$x$(gam)}; 
			\end{tikzpicture}
	\end{center}
	Hãy ước lượng số trung bình và mốt của mẫu số liệu  ghép nhóm trên (kết quả làm tròn đến hàng phần trăm).
	 \loigiai{
		\begin{center}
			\begin{tabular}{|l|c|c|c|c|c|}
				\hline Khối lượng & {$[70;80)$}  & {$ [80;90)$} & {$[90;100) $} & {$[100;110)$}& {$[110;120)$}\\
				\hline Gía trị đại diện &  $75$ & $85$ & $95$ & $105$& $115$\\
				\hline Số củ & $3$ & $5$ &  $12$ & $6$& $3$\\
				 \hline
			\end{tabular}
		\end{center}
		Khối lượng trung bình mỗi củ khoai tây là  $$\overline{x}=\dfrac{75\cdot 3+85\cdot 5+95\cdot 12+105\cdot 6+115\cdot 3}{30}=\dfrac{553}{6}\approx 92{,}17\text{(g)}.$$
		Nhóm chứa mốt của  mẫu số liệu là $[90;100)$.\\
		Do  đó $u_m=90;n_{m-1}=5;n_{m+1}=6;u_{m+1}-u_m=100-90=10$.\\
		mốt của  mẫu số liệu ghép nhóm là $M_o=90+\dfrac{12-5}{(12-5)+(12-6)}\cdot 10=\dfrac{1240}{13}\approx 95{,}38\text{(g)}$.
	} 
\end{ex}

% nguồn: Ok_TuanLuan_Toan11_Lan2
\begin{ex}%[1D5H1-4]%[ma:3cb92ffe]
	Hãy sử dụng dữ liệu ở để tư vấn cho đại lí bảo hiểm xác định khách hàng nam và nữ ở tuổi  nào hay mua bảo hiểm nhất.
	Số khách hàng mua bảo  hiểm ở từng độ tuổi được thống kê như sau:
	\begin{center}
		\begin{tabular}{|c|c|c|c|c|c|}
			\hline Độ tuổi & $[20;30)$ & $[30;40)$  & $[40;50)$ & $[50;60)$ & $[60;70)$\\
			 \hline
			Số khách hàng nam & $4$  & $6$ & $10$ & $7$ & $3$ \\
			\hline 
			Số khách hàng nam  & $3$ & $9$ & $6$ & $4$ & $2$ \\
			\hline 
		\end{tabular}
	\end{center}
	\loigiai{ 
		Nhóm chứa mốt của mẫu số liệu khách hàng nam  là $\left[40;50\right)$.\\
		Do đó $u_m=40$, $n_{m-1}=6$, $n_{m+1}=7$,  $u_{m+1}-u_m=50-40=10$.\\
		Mốt của mẫu số  liệu nhóm khách hàng nam là
		 \[M_o=40+\dfrac{10-6}{\left(10-6\right)+\left(10-7\right)} \cdot 10=45{,}7.
		\] 
		Dựa vào kết quả trên ta có thể dự đoán được khách hàng nam $46$ tuổi có nhu cầu mua bảo hiểm cao nhất. Nhóm chứa mốt của mẫu số liệu khách hàng  nữ là $\left[30;40\right)$.\\
		Do đó $u_m=30$, $n_{m-1}=3$,  $n_{m+1}=6$, $u_{m+1}-u_m=40-30=10$.\\
		Mốt  của mẫu số liệu nhóm khách hàng nam là
		\[M_o=30+\dfrac{9-3}{\left(9-3\right)+\left(9-6\right)}  \cdot 10=36{,}7.
		 \]
		Dựa vào kết quả trên ta có thể dự đoán được khách hàng nữ $37$ tuổi có nhu cầu mua bảo  hiểm cao nhất.
	} 
\end{ex}

% nguồn: Ok_TuanLuan_Toan11_Lan2
\begin{ex}%[1D5H1-4]%[ma:42507cd2]
	 Mức thưởng tết (triệu đồng) mà các công nhân một nhà máy nhận được như sau:
	\begin{center}
		\begin{tabular}{|*5{c|}}
			\hline 
			 Mức thưởng & $[5; 10)$& $[10; 15)$& $[15; 20)$ & $[20; 25]$\\
			 \hline
			Số công  nhân & 13 & 35 & 47 & 25 \\
			 \hline
		\end{tabular}
	\end{center}
	 Tìm mốt của mẫu số liệu ghép nhóm. Cho biết ý nghĩa của giá trị thu được.
	 \loigiai{
		Cỡ mẫu $n=13+35+47+25=120$. Số  công nhân có mức thưởng tết từ $15$ đến dưới $20$ triệu đồng là nhiều nhất nên nhóm chứa mốt là nhóm $\left[15; 20 \right)$.\\
		Ta có $a_j=15$; $m_j=47$; $m_{j-1}=35$; $m_{j+1}=25$; $h=5$. Do đó, mốt của mẫu số liệu là  \[M_o=15+\dfrac{(47-35)}{(47-35)+(47-25)}\cdot 5\approx 16{,}76.\]
		Ý nghĩa: Số công nhân nhận được mức thưởng tết  khoảng $16{,}76$ triệu đồng là cao nhất.
	} 
\end{ex}

% nguồn: Ok_TuanLuan_Toan11_Lan2
\begin{ex}%[1D5H1-4]%[ma:c971b6ae]
	Thời gian (phút) để học  sinh hoàn thành một câu hỏi thi được cho như sau:
	\begin{center}
		\begin{tabular}{|c|c|c|c|c|c|}
			\hline 
			Thời gian  (phút) & $[0{,}5; 10{,}5)$& $[10{,}5; 20{,}5)$& $[20{,}5; 30{,}5)$&$[30{,}5; 40{,}5)$&$[40{,}5; 50{,}5)$ \\
			\hline 
			Số học sinh & $2$ & $10$ & $6$  & $4$ & $3$ \\
			 \hline
		\end{tabular}
	\end{center}
	Tìm mốt  của mẫu số liệu ghép nhóm này.
	 \loigiai{
		 Tần số lớn nhất là $10$ nên nhóm chứa mốt là nhóm $\left[10{,}5;20{,}5 \right)$.\\
		Ta có, $j=2$; $a_2=10{,}5$; $m_2=10$; $m_1=2$; $m_3=6$; $h=20{,}5-10{,}5=10$. Do đó, mốt của mẫu số liệu ghép nhóm này là $M_o=10{,}5+\dfrac{10-2}{(10-2)+(10-6)}\cdot 10=\dfrac{103}{6}\approx  17{,}17$.
	} 
\end{ex}

% nguồn: Ok_TuanLuan_Toan11_Lan2
\begin{ex}%[1D5H1-4]%[ma:8f59b52b]
	Một thư viện thống kê số  lượng sách được mượn mỗi ngày trong ba tháng ở bảng sau
	\begin{center}
		\begin{tabular}{|c|c|c|c|c|c|c|c|}
			 \hline
			Số sách & $[16;20]$ & $[21;25]$ & $[26;30]$ & $[31;35]$  & $[36;40]$ & $[41;45]$ & $[46;50]$\\
			\hline 
			Số ngày & $3$ & $6$ &  $15$ & $27$ & $22$ & $14$ & $5$\\
			\hline 
		\end{tabular}
	\end{center}
	Hãy ước lượng số trung bình và  mốt của mẫu số liệu ghép nhóm trên.
	 \loigiai{
		 Số liệu trên được hiệu chỉnh như sau
		\begin{center}
			 \resizebox{\textwidth}{!}{
			\begin{tabular}{|c|c|c|c|c|c|c|c|}
				\hline 
				Số sách & $[15{,}5;20{,}5)$ & $[20{,}5;25{,}5)$ & $[25{,}5;30{,}5)$  & $[30{,}5;35{,}5)$ & $[35{,}5;40{,}5)$ & $[40{,}5;45{,}5)$ & $[45{,}5;50{,}5)$\\
				 \hline
				Giá trị  đại diện & $18$ & $23$ & $28$ & $33$ & $38$ & $43$ & $48$\\
				\hline 
				Số  ngày & $3$ & $6$ & $15$ & $27$ & $22$ & $14$ & $5$\\
				 \hline
			\end{tabular}
			} 
		\end{center}
		Số sách  được mượn trung bình mỗi ngày xấp xỉ bằng
		\[\dfrac{18\cdot 3+23\cdot  6+28\cdot 15+33\cdot 27+38\cdot 33+43\cdot 14+48\cdot 5}{92}=34{,}6.\]
		Nhóm  chứa mốt của mẫu số liệu là $\left[30{,}5;35{,}5\right)$.\\
		Do đó $u_m=30{,}5$;  $n_{m-1}=15$, $n_{m+1}=22$, $u_{m+1}-u_m=35{,}5-30{,}5=5$.\\
		Mốt của  mẫu số liệu ghép nhóm là
		 $$M_0=30{,}5+\dfrac{27-15}{(27-15)+(27-22)}\cdot 5=34.$$
	} 
\end{ex}

% nguồn: Ok_TuanLuan_Toan11_Lan2
\begin{ex}%[1D5H1-4]%[ma:14dc3047]
	Kết quả đo chiều cao của $200$ cây keo $3$ năm tuổi ở một nông trường được  biểu diễn ở biểu đồ dưới đây.
	\begin{center}
		\begin{tikzpicture}[>=stealth,line join=round,line cap=round,font=\footnotesize,scale=0.7,xscale=0.6]
				\foreach \y in {0,10,...,70}{ 
					\node[left] at  (0,\y/10){$\y$};
				} 
				\foreach  \x/\y/\a in {0/20/$[8{,}5; 8{,}8)$,6/35/$[8{,}8; 9{,}1)$,12/60/$[9{,}1; 9{,}4)$,18/55/$[10; 15)$,24/30/$[9{,}7; 10)$}{
					\fill[blue!15] (\x,0) rectangle (\x+6,\y/10); 
					\node[above] at (\x+3,\y/10){$\y$}; 
					 \node[below] at (\x+3,0){\a};
					 \draw (\x,0) rectangle (\x+6,\y/10);
				} 
				 \draw[<->] (0,8) node[above left]{Số cây}--(0,0)--(32,0) node[below right]{(m)};
				 \node at (13,9){\bfseries Chiều cao của $200$ cây keo $3$ năm tuổi};
			\end{tikzpicture}
	\end{center}
	Hãy  ước lượng số trung bình và mốt của mẫu số liệu ghép nhóm trên.
	\loigiai{ 
		Chiều cao của  $200$ cây keo được thống kê như bảng sau
		\begin{center}
			\begin{tabular}{|c|c|c|c|c|c|}
				 \hline
				Chiều cao &  $[8{,}5;8{,}8)$ & $[8{,}8;9{,}1)$ & $[9{,}1;9{,}4)$ & $[9{,}4;9{,}7)$ & $[9{,}7;10)$ \\
				 \hline
				Giá trị đại diện & $8{,}65$ & $8{,}95$ & $9{,}25$ & $9{,}55$  & $9{,}85$ \\
				\hline 
				Số cây & $20$ & $35$ & $60$ & $55$ &  $30$ \\
				 \hline
			\end{tabular}
		\end{center}
		Chiều cao trung bình của $200$ cây xấp xỉ bằng 
		\[\dfrac{8{,}65\cdot 20+8{,}95\cdot 35+9{,}25\cdot 60+9{,}55\cdot  55+9{,}85\cdot 30}{200}=9{,}31\text{(m)}.\]
		 Nhóm chứa mốt của mẫu số liệu là $\left[9{,}1;9{,}4\right)$.\\
		Do đó $u_m=9{,}1$; $n_{m-1}=35$;  $n_{m+1}=55$; $u_{m+1}-u_m=9{,}4-9{,}1=0{,}3$.\\
		 Mốt của mẫu số liệu trên là
		$$M_0=9{,}1+\dfrac{60-35}{(60-35)+(60-55)}\cdot  0{,}3=9{,}35\text{ (m)}.$$
	} 
\end{ex}

% nguồn: Ok_TuanLuan_Toan11_Lan2
\begin{ex}%[1D5H1-4]%[ma:172f0da6]
	Một công ty xây dựng khảo sát khách hàng xem họ có nhu cầu mua  nhà ở mức giả nào. Kết quả khảo sát được ghi lại ở bảng sau:
	\begin{center}
		\begin{tabular}{|c|c|c|c|c|c|}
			\hline Mức giá (triệu đồng/$\mathrm{m}^2$) & $[10;14)$ & $[14;18)$ & $[18;22)$ & $[22;26)$ &  $[26;30)$\\
			 \hline
			Số khách hàng &  $54$ & $78$ & $120$ & $45$ & $12$ \\
			\hline 
		\end{tabular}
	\end{center}
	\begin{enumerate}[a)]
		\item Tìm mốt của mẫu số liệu ghép nhóm  trên.
		\item Công ty nên xây nhà ở mức giá nào để nhiều  người có nhu cầu mua nhất?
	\end{enumerate}
	\loigiai{ 
		\begin{enumerate}
			\item Nhóm chứa mốt của  mẫu số liệu trên là nhóm $\left[18;22\right)$.\\
			Do  đó $u_m=18$, $n_{m-1}=78$, $n_m=120$, $n_{m+1}=45$, $u_{m+1}-u_m=22-18=4$.\\
			Mốt của mẫu số liệu ghép nhóm là 
			\[M_o=18+\dfrac{120-78}{\left(120-78\right)+\left(120-45\right)} \cdot  4=\dfrac{758}{39} \approx 19{,}4.
			 \]
			\item Dựa vào kết quả trên ta có thể dự đoán rằng nếu công ty xây nhà ở mức giá $19{,}4$ triệu đồng/ $\mathrm{m}^2$ thì sẽ có nhiều người có nhu cầu mua  nhất.
		\end{enumerate}
	 }
\end{ex}

% nguồn: Ok_TuanLuan_Toan11_Lan2
\begin{ex}%[1D5H1-4]%[ma:a30d768b]
	Tìm hiểu thời gian xem ti vi trong tuần trước  (đơn vị: giờ) của một số học sinh thu được kết quả sau:
	\begin{center}
		\begin{tabular}{|c|c|c|c|c|c|}
			\hline 
			Thời gian (giờ) & $[0; 5)$ &$[5; 10)$ & $[10; 15)$&$[15;  20)$&$[20; 25)$ \\
			 \hline
			Số học sinh &$8$ &$16$ &$4$ &$2$ &$2$ \\ 
			\hline 
		\end{tabular}
	\end{center}
	\begin{listEX}
		\item Có thể tìm được giá trị chính xác cho mốt của  mẫu số liệu gốc về thời gian xem ti vi của học sinh không?
		\item Mốt thuộc nhóm nào là hợp lí nhất? Nên lấy số nào trong nhóm để ước  lượng cho mốt?
	\end{listEX}
	 \loigiai{
		\begin{listEX}
			\item Không thể tính được giá trị chính xác cho mốt của mẫu số liệu gốc về thời gian xem ti vi của học  sinh, do không có thời gian cụ thể của từng học sinh.
			\item Tần số lớn nhất là 16 nên mốt thuộc nhóm $\left[5;10 \right)$ là  hợp lí nhất. Ta ước lượng mốt của mẫu số liệu bằng cách xác định số thứ tự của nhóm chứa mốt là $j=2$; $a_j=a_2=5$; $m_2=16$; $m_1=8$; $m_3=4$; độ dài của nhóm $h=5$.\\
			Do đó, mốt của mẫu  số liệu xấp xỉ bằng $5+\dfrac{16-8}{(16-8)+(16-4)} \cdot 5=7$.
	\end{listEX}}
\end{ex}
